\documentclass[10pt]{scrartcl}
\usepackage[mdthm,colorsec,noauthor,noasy]{evan}
\usepackage{handout}
\newcommand{\ip}[2]{\langle #1,#2\rangle}
\DeclareMathOperator{\tr}{tr}\DeclareMathOperator{\nul}{null}\DeclareMathOperator{\range}{range}\DeclareMathOperator{\spn}{span}
\newcommand{\LL}{\mathcal L}\newcommand{\PP}{\mathcal P}\newcommand{\MM}{\mathcal M}
\title{Math 110 --- Week 11}
\subtitle{Adjoints, self-adjoint and normal operators, the spectral theorem (Notes Ch.~19--20, LADR 7A--7B, HW11)}
\author{}\date{}\ihead{}
\begin{document}
\maketitle

\section{Facts used from earlier weeks}

Throughout, $V$ and $W$ are finite-dimensional inner product spaces over $\FF=\RR$ or $\CC$, with $n=\dim V$ and $m=\dim W$. A bar $\overline{z}$ denotes complex conjugation; for real scalars it changes nothing.

\begin{enumerate}[label=(T\arabic*),leftmargin=2.6em,itemsep=8pt]
  \ii \target{T1}\tool{Rules of the inner product (\xweek{math110}{2026-10-26-trace-inner-products}{Week~9}).} For all $u,v,w\in V$ and $\lambda\in\FF$: linearity in the first slot $\langle \lambda u+v,w\rangle = \lambda\langle u,w\rangle + \langle v,w\rangle$; conjugate-linearity in the second slot $\langle u,\lambda v+w\rangle = \overline{\lambda}\langle u,v\rangle + \langle u,w\rangle$; conjugate symmetry $\langle u,v\rangle = \overline{\langle v,u\rangle}$; and \textbf{positive definiteness}: $\langle v,v\rangle \ge 0$ for every $v\in V$ (positivity), and $\langle v,v\rangle = 0 \iff v=0$ (definiteness). The norm is $\|v\|=\sqrt{\langle v,v\rangle}$, so $\|v\|^2=\langle v,v\rangle$.

  \ii \target{T2}\tool{Cauchy--Schwarz (\xlink{math110}{2026-10-26-trace-inner-products}{prop-82}{Week~9, Prop.~82}).} $|\langle u,v\rangle| \le \|u\|\,\|v\|$, with equality if and only if one of $u,v$ is a scalar multiple of the other.

  \ii \target{T3}\tool{Orthonormal bases (\xlink{math110}{2026-11-02-orthogonality-projections}{prop-84}{Week~10, Props.~84--87}).} In an orthonormal basis $e_1,\dots,e_n$ of $V$, every vector expands as $v = \langle v,e_1\rangle e_1 + \dots + \langle v,e_n\rangle e_n$, and $\|a_1e_1+\dots+a_ne_n\|^2 = |a_1|^2+\dots+|a_n|^2$. Every orthonormal list is linearly independent, so $n$ orthonormal vectors in $V$ form a basis. Here $\delta_{ij}=1$ if $i=j$ and $0$ otherwise. Gram--Schmidt (\xlink{math110}{2026-11-02-orthogonality-projections}{prop-87}{Prop.~87}) turns a basis $v_1,\dots,v_k$ into an orthonormal basis $e_1,\dots,e_k$ with $\spn(e_1,\dots,e_j)=\spn(v_1,\dots,v_j)$ for every $j$; applied to any subspace, it yields an orthonormal basis of that subspace.

  \ii \target{T4}\tool{Riesz representation (\xlink{math110}{2026-11-02-orthogonality-projections}{prop-88}{Week~10, Prop.~88}).} For every linear functional $\psi\in V'$ there exists a unique $u\in V$ such that $\psi(v)=\langle v,u\rangle$ for all $v\in V$. In particular, if $\langle v,u\rangle=\langle v,u'\rangle$ for every $v\in V$, then $u=u'$ (set $v=u-u'$ to obtain $\|u-u'\|^2=0$).

  \ii \target{T5}\tool{Orthogonal complements (\xlink{math110}{2026-11-02-orthogonality-projections}{prop-89}{Week~10, Props.~89--90}).} For any subspace $U\subseteq V$: $V=U\oplus U^\perp$, $\dim U^\perp = n-\dim U$, and $(U^\perp)^\perp = U$.

  \ii \target{T6}\tool{Dual map (\xweek{math110}{2026-09-28-quotients-duality}{Week~5}).} For $T\in\LL(V,W)$, its dual map $T'\in\LL(W',V')$ is defined by $T'(\varphi)=\varphi\circ T$.

  \ii \target{T7}\tool{Eigenvalues (\xlink{math110}{2026-10-05-eigenvalues-pagerank}{prop-53}{Week~6, Props.~53--54}).} On a nonzero finite-dimensional complex space every operator has an eigenvalue. Eigenvectors corresponding to distinct eigenvalues are linearly independent.

  \ii \target{T8}\tool{Minimal polynomial and diagonalizability (\xweek{math110}{2026-10-12-diagonalization-minimal-polynomial}{Week~7}).} $m_T$ is the monic polynomial of smallest degree with $m_T(T)=0$; $q(T)=0 \iff m_T\mid q$ (\xlink{math110}{2026-10-12-diagonalization-minimal-polynomial}{prop-65}{Prop.~65}); the eigenvalues of $T$ are the roots of $m_T$ in $\FF$ (\xlink{math110}{2026-10-12-diagonalization-minimal-polynomial}{prop-66}{Prop.~66}); $T$ is diagonalizable $\iff$ $V$ has an eigenbasis $\iff V=E(\lambda_1,T)\oplus\dots\oplus E(\lambda_m,T)\iff\sum_i\dim E(\lambda_i,T)=n$ (\xlink{math110}{2026-10-12-diagonalization-minimal-polynomial}{prop-62}{Prop.~62}) $\iff$ $m_T$ is a product of distinct linear factors (\xlink{math110}{2026-10-12-diagonalization-minimal-polynomial}{prop-68}{Prop.~68}). If $T$ is diagonalizable and $U$ is $T$-invariant, then $T|_U$ is diagonalizable (Week~7, F3).

  \ii \target{T9}\tool{Real polynomials (\xlink{math110}{2026-10-05-eigenvalues-pagerank}{name-fundamental-theorem-of-algebra}{Week~6}; cited in Week~7, F11).} Every nonconstant polynomial in $\PP(\RR)$ factors into real linear factors and real irreducible quadratics $x^2+bx+c$ with $b^2<4c$.

  \ii \target{T10}\tool{Upper-triangular matrices (\xlink{math110}{2026-10-12-diagonalization-minimal-polynomial}{prop-59}{Week~7, Props.~59, 61}).} In a basis $v_1,\dots,v_n$, $\MM(T)$ is upper-triangular $\iff$ $\spn(v_1,\dots,v_k)$ is $T$-invariant for each $k$. Over $\CC$, such a basis always exists.

  \ii \target{T11}\tool{Jordan chains (\xlink{math110}{2026-10-19-generalized-eigenvectors-jordan}{prop-76}{Week~8, Prop.~76}).} Over $\CC$, $V$ has a basis of Jordan chains: lists $v_1,\dots,v_k$ with $Tv_1=\lambda v_1$ and $Tv_j=\lambda v_j+v_{j-1}$ ($j\ge2$). $T$ is diagonalizable $\iff$ every Jordan chain has length $1$.

  \ii \target{T12}\tool{$2\times2$ matrices (\xweek{math110}{2026-10-05-eigenvalues-pagerank}{Week~6}).} The characteristic polynomial of a $2\times2$ matrix $A$ is $x^2-(\tr A)x+\det A$, whose roots are its eigenvalues.

  \ii \target{T13}\tool{Rank--nullity (\xweek{math110}{2026-09-14-linear-maps-matrices}{Week~3}).} For $T\in\LL(V,W)$, $\dim\nul T+\dim\range T=\dim V$; an injective operator on a finite-dimensional space is invertible.

  \ii \target{T14}\tool{Projectors (\xlink{math110}{2026-11-02-orthogonality-projections}{prop-91}{Week~10, Prop.~91, Def.~48}).} If $P^2=P$, then $V=\nul P\oplus\range P$ and $P|_{\range P}=\operatorname{id}$. The orthogonal projection $P_U$ onto $U$ acts as the identity on $U$ and $0$ on $U^\perp$.

  \ii \target{T15}\tool{Bases of a direct sum (\xweek{math110}{2026-10-19-generalized-eigenvectors-jordan}{Week~8}).} If $V=U_1\oplus\dots\oplus U_m$, concatenating bases of $U_1,\dots,U_m$ produces a basis of $V$.
\end{enumerate}

\section{This week}

\subsection{The adjoint (Ch.~19)}

\begin{qbox}
  On $\PP_1(\RR)$ with the inner product $\langle f,g\rangle=\int_0^1 f(x)g(x)\,dx$, let $D\colon\PP_1(\RR)\to\PP_1(\RR)$ be differentiation. Does there exist an operator $D^*$ satisfying $\langle Df,g\rangle=\langle f,D^*g\rangle$ for all $f,g$, and what is $D^*(1)$? In the basis $1,x$, the matrix of $D$ is
  \[
    \MM(D)=\begin{bmatrix}0&1\\0&0\end{bmatrix}.
  \]
  Is $D^*$ represented by its transpose $\MM(D)^{\mathsf T}$?
\end{qbox}

In an inner product space we measure geometric alignment via $\langle Tv,w\rangle$. Moving $T$ onto the other vector $w$ defines a companion map $T^*\colon W\to V$ that realizes all the linear-algebraic duality of Week~5 directly inside the inner product spaces $V$ and $W$.

\begin{definition*}[Def.~49 / LADR 7A: Adjoint]\target{def-49}\target{name-adjoint}
  Let $T\in\LL(V,W)$. The \textbf{adjoint} of $T$ is the map $T^*\colon W\to V$ defined by
  \[
    \langle Tv,w\rangle = \langle v,T^*w\rangle \qquad \text{for all } v\in V,\; w\in W.
  \]
  (Many physics and engineering texts write $T^\dagger$.)
\end{definition*}

\begin{ideabox}
  \textbf{Why $T^*$ exists and is unique, and how it relates to the dual map $T'$.} Fix $w\in W$. Because the inner product is linear in the first slot \ilink{T1}{(T1)}, the scalar function $\varphi_w(v)=\langle Tv,w\rangle$ is a linear functional in $V'$. By the Riesz representation theorem \ilink{T4}{(T4)}, there is a \emph{unique} vector $u\in V$ with $\langle Tv,w\rangle=\langle v,u\rangle$ for every $v\in V$; we set $T^*w=u$.

  Writing $\Gamma_V(v)=\langle\cdot,v\rangle\in V'$ and $\Gamma_W(w)=\langle\cdot,w\rangle\in W'$ for the Riesz identifications from Week~10, the defining identity reads
  \[
    \langle \cdot,T^*w\rangle = \langle T(\cdot),w\rangle = \langle\cdot,w\rangle\circ T = T'(\Gamma_W(w)),
  \]
  i.e.\ $\Gamma_V(T^*w)=T'(\Gamma_W(w))$. Thus $T^*=\Gamma_V^{-1}\circ T'\circ \Gamma_W$: \emph{the adjoint is the dual map $T'$ transported back from $W'\to V'$ to $W\to V$ via the Riesz inner product}. Because $\Gamma_V$ depends on the chosen inner product and satisfies $\Gamma_V(\lambda v)=\overline{\lambda}\,\Gamma_V(v)$ over $\CC$, the adjoint $T^*$ depends on the inner product and conjugates scalars.
\end{ideabox}

\begin{fig}
\[
\begin{tikzcd}[
  arrow style=tikz,
  diagrams={>={Stealth[length=4.5pt,width=2.8pt]}},
  column sep=8.6em,
  row sep=3.8em,
  every label/.append style={outer sep=2.5pt}
]
  W \arrow[r, draw=figblue, thick, "{\color{figblue!90!black} T^* = \Gamma_V^{-1}\circ T'\circ\Gamma_W}"]
    \arrow[d, draw=figgreen!85!black, thick, "{\color{figgreen!90!black}\substack{\Gamma_W \\ \scriptstyle w\mapsto\langle\cdot,w\rangle}}"']
  & V \arrow[d, draw=figgreen!85!black, thick, "{\color{figgreen!90!black}\substack{\Gamma_V \\ \scriptstyle v\mapsto\langle\cdot,v\rangle}}"] \\[2pt]
  W' \arrow[r, draw=figslate, thick, "{\color{black} T'\colon \varphi\mapsto \varphi\circ T}"']
  & V'
\end{tikzcd}
\]
\figcap{Both paths from $W$ to $V'$ agree: $\Gamma_V(T^*w)=T'(\Gamma_W(w))$. Because the Riesz maps $\Gamma_W,\Gamma_V$ (green) are conjugate-linear over $\CC$, transporting the dual map $T'$ back to $T^*\colon W\to V$ (blue) introduces complex conjugation on scalars.}
\end{fig}

\begin{example*}[Ex.~72: Adjoint of differentiation by solving equations]\target{ex-72}
  In the question, write $D^*(1)=ax+b$. Applying $\langle f,D^*(1)\rangle = \langle Df,1\rangle$ to the basis functions $f=1$ and $f=x$ gives two equations:
  \[
    \langle 1,ax+b\rangle = \langle D(1),1\rangle = 0, \qquad
    \langle x,ax+b\rangle = \langle D(x),1\rangle = \langle 1,1\rangle = 1.
  \]
  Computing the integrals gives $\langle 1,ax+b\rangle=\frac12 a+b=0$ and $\langle x,ax+b\rangle=\frac13 a+\frac12 b=1$. Solving yields $a=12$, $b=-6$, so $D^*(1)=12x-6$.

  Similarly, writing $D^*(x)=cx+d$ gives $\frac12 c+d=\langle D(1),x\rangle=0$ and $\frac13 c+\frac12 d=\langle D(x),x\rangle=\langle 1,x\rangle=\frac12$, whence $c=6$, $d=-3$ and $D^*(x)=6x-3$.
\end{example*}

\begin{proposition*}[Prop.~93 / LADR 7.4, 7.9: Linearity and conjugate-transpose matrix]\target{prop-93}
  If $T\in\LL(V,W)$, then $T^*\in\LL(W,V)$ is linear. Moreover, if $e_1,\dots,e_n$ is an orthonormal basis of $V$, $f_1,\dots,f_m$ is an orthonormal basis of $W$, and $\MM(T)=(a_{ij})$ with respect to these bases, then
  \[
    \MM(T^*) = \bigl(\overline{a_{ji}}\bigr) = \MM(T)^*,
  \]
  the conjugate transpose of $\MM(T)$.
\end{proposition*}
\begin{ideabox}
  Both statements come from testing inner products against arbitrary vectors $v\in V$. To prove $T^*$ is linear, we show $T^*(\lambda w_1+w_2)$ and $\lambda T^*w_1+T^*w_2$ have the exact same inner product with every $v\in V$, then invoke uniqueness \ilink{T4}{(T4)}. To find the $(k,j)$-entry of $\MM(T^*)$, remember that in an orthonormal basis $e_1,\dots,e_n$, coordinates are inner products \ilink{T3}{(T3)}: the $e_k$-coordinate of $T^*f_j$ is $\langle T^*f_j,e_k\rangle$, and moving $T^*$ across flips it into $\overline{a_{jk}}$.
\end{ideabox}
\begin{proof}
  \begin{steps}
    \ii \emph{Linearity.} For $v\in V$, $w_1,w_2\in W$, and $\lambda\in\FF$, using \ilink{T1}{(T1)} and Def.~49,
    \[
      \langle v, T^*(\lambda w_1+w_2)\rangle
      = \langle Tv,\lambda w_1+w_2\rangle
      = \overline{\lambda}\langle Tv,w_1\rangle + \langle Tv,w_2\rangle
      = \overline{\lambda}\langle v,T^*w_1\rangle + \langle v,T^*w_2\rangle
      = \langle v,\lambda T^*w_1+T^*w_2\rangle.
    \]
    Since this holds for all $v\in V$, \ilink{T4}{(T4)} gives $T^*(\lambda w_1+w_2)=\lambda T^*w_1+T^*w_2$.
    \ii \emph{Matrix.} Write $\MM(T^*)=(b_{kj})$. Column $j$ holds the $e$-coordinates of $T^*f_j$, so by \ilink{T3}{(T3)} $b_{kj}=\langle T^*f_j,e_k\rangle$. Moving $T^*$ to the other slot via conjugate symmetry \ilink{T1}{(T1)},
    \[
      b_{kj} = \langle T^*f_j,e_k\rangle = \overline{\langle e_k,T^*f_j\rangle} = \overline{\langle Te_k,f_j\rangle} = \overline{a_{jk}},
    \]
    because $a_{jk}=\langle Te_k,f_j\rangle$ is the $f_j$-coordinate of $Te_k$ \ilink{T3}{(T3)}.
  \end{steps}
\end{proof}

By linearity of $D^*$ in Example~72, $D^*(ax+b)=a(6x-3)+b(12x-6)=(3a+6b)(2x-1)$.

\begin{example*}[Ex.~73: Example~72 via an orthonormal basis]\target{ex-73}
  Orthonormalizing $\{1,x\}$ on $\PP_1(\RR)$ gives $e_1=1$ and $e_2=\sqrt{3}(2x-1)$ (check $\int_0^1 3(2x-1)^2\,dx=1$ and $\int_0^1 (2x-1)\,dx=0$). Since $D(e_1)=0$ and $D(e_2)=2\sqrt{3}=2\sqrt{3}\,e_1$, Prop.~93 yields
  \[
    \MM(D)=\begin{bmatrix}0&2\sqrt{3}\\0&0\end{bmatrix}, \qquad
    \MM(D^*)=\begin{bmatrix}0&0\\2\sqrt{3}&0\end{bmatrix}.
  \]
  Hence $D^*(e_1)=2\sqrt{3}\,e_2=12x-6$ and $D^*(e_2)=0$. Writing $x=\frac{1}{2\sqrt{3}}e_2+\frac12 e_1$ recovers $D^*(x)=\frac12 D^*(1)=6x-3$.
\end{example*}

\noindent\textbf{The question answered.} $D^*$ exists and $D^*(1)=12x-6$. Taking the transpose of $\MM(D)$ in $\{1,x\}$ falsely gives $D^*(1)=x$: Prop.~93 requires an \emph{orthonormal} basis, whereas $\langle 1,x\rangle=\frac12\neq 0$.

\begin{proposition*}[LADR 7.5: Algebraic rules for adjoints]\target{name-ladr-7-5}
  For $S,T\in\LL(V,W)$, $\lambda\in\FF$, and $R\in\LL(U,V)$ ($U$ an inner product space):
  \[
    (S+T)^*=S^*+T^*,\qquad (\lambda T)^*=\overline{\lambda}\,T^*,\qquad (T^*)^*=T,\qquad I^*=I,\qquad (TR)^*=R^*T^*.
  \]
  In particular $(T-\lambda I)^*=T^*-\overline{\lambda}\,I$; and if $T$ is invertible, then $T^*$ is invertible with $(T^*)^{-1}=(T^{-1})^*$.
\end{proposition*}
\begin{proof}
  Each identity follows by computing $\langle\text{(map)}v,w\rangle$ for arbitrary vectors and using uniqueness in \ilink{T4}{(T4)}:
  \begin{steps}
    \ii $(S+T)^*$: $\langle(S+T)v,w\rangle = \langle Sv,w\rangle+\langle Tv,w\rangle = \langle v,S^*w\rangle+\langle v,T^*w\rangle = \langle v,(S^*+T^*)w\rangle$.
    \ii $(\lambda T)^*$: $\langle\lambda Tv,w\rangle = \lambda\langle v,T^*w\rangle = \langle v,\overline{\lambda}\,T^*w\rangle$ (conjugate-linear second slot, T1).
    \ii $(T^*)^*$: $\langle T^*w,v\rangle = \overline{\langle v,T^*w\rangle} = \overline{\langle Tv,w\rangle} = \langle w,Tv\rangle$ for all $w\in W,v\in V$.
    \ii $I^*$: $\langle Iv,w\rangle = \langle v,w\rangle = \langle v,Iw\rangle$.
    \ii $(TR)^*$: $\langle TRu,w\rangle = \langle Ru,T^*w\rangle = \langle u,R^*T^*w\rangle$ (order reverses).
    \ii Shifts \& inverses: combine (1), (2), (4) for $(T-\lambda I)^*=T^*-\overline{\lambda}I$; applying (5) to $TT^{-1}=I$ and $T^{-1}T=I$ yields $(T^{-1})^*T^*=I$ and $T^*(T^{-1})^*=I$.
  \end{steps}
\end{proof}

\begin{theorem*}[Prop.~94 / LADR 7.6: Fundamental theorem of adjoint subspaces]\target{prop-94}
  For every $T\in\LL(V,W)$:
  \begin{enumerate}
    \ii[(i)] $\nul T^* = (\range T)^\perp$;
    \ii[(ii)] $\range T^* = (\nul T)^\perp$.
  \end{enumerate}
  Consequently, $V = \nul T \oplus \range T^*$ and $W = \range T \oplus \nul T^*$.
\end{theorem*}
\begin{ideabox}
  A vector $w\in W$ lies in $\nul T^*$ precisely when $T^*w=0$, which by nondegeneracy \ilink{T4}{(T4)} is equivalent to $\langle v,T^*w\rangle=0$ for all $v\in V$. Rewriting $\langle v,T^*w\rangle$ as $\langle Tv,w\rangle$ characterizes $T^*w=0$ as ``$w$ is orthogonal to every output $Tv\in\range T$'', yielding $\nul T^*=(\range T)^\perp$. Swapping $T$ and $T^*$ (since $(T^*)^*=T$) and taking $(\cdot)^\perp$ gives (ii). Together with $U\oplus U^\perp$, this splits both domain $V$ and codomain $W$ orthogonally.
\end{ideabox}
\begin{proof}
  \begin{steps}
    \ii \textbf{(i).} $w\in\nul T^* \iff T^*w=0 \iff \langle v,T^*w\rangle=0\ \forall v\in V$ \ilink{T4}{(T4)} $\iff \langle Tv,w\rangle=0\ \forall v\in V \iff w\in(\range T)^\perp$.
    \ii \textbf{(ii).} Apply (i) to $T^*\in\LL(W,V)$ using $(T^*)^*=T$: $\nul T = (\range T^*)^\perp$. Taking orthogonal complements of both sides and using $(U^\perp)^\perp=U$ \ilink{T5}{(T5)} yields $(\nul T)^\perp = \range T^*$.
    \ii The direct sums follow immediately from \ilink{T5}{(T5)}: $V=\nul T\oplus(\nul T)^\perp = \nul T\oplus\range T^*$ and $W=\range T\oplus(\range T)^\perp=\range T\oplus\nul T^*$.
  \end{steps}
\end{proof}

\begin{fig}
\begin{tikzpicture}[scale=1.0,>=Stealth]
  \begin{scope}
    \node[font=\footnotesize,align=center,anchor=south] at (1.15,3.88) {\textbf{Domain $V$}\\[-1pt]\scriptsize $\range T^*\oplus\nul T$};
    \node[font=\footnotesize,align=center,anchor=south] at (6.55,3.88) {\textbf{Codomain $W$}\\[-1pt]\scriptsize $\range T\oplus\nul T^*$};

    % Domain V: range T* (blue) and null T (red)
    \draw[draw=figblue!85!black,line width=0.65pt,fill=figblue!6] (0,2.50) rectangle (2.30,3.75);
    \node[font=\footnotesize,align=center,text=figblue!90!black] at (1.15,3.13) {$\range T^*$\\[-1pt]\scriptsize $=(\nul T)^\perp$};

    \node[font=\scriptsize] at (0.95,1.66) {$\perp$};
    \fill[black] (2.30,1.66) circle (1.8pt) node[left=3pt,font=\scriptsize] {$0_V$};

    \draw[draw=figred!85!black,dashed,line width=0.6pt,fill=figred!5] (0,0.0) rectangle (2.30,1.25);
    \node[font=\footnotesize,text=figred!90!black] at (1.15,0.62) {$\nul T$};

    % Codomain W: range T (blue) and null T* (red)
    \draw[draw=figblue!85!black,line width=0.65pt,fill=figblue!6] (5.40,2.50) rectangle (7.70,3.75);
    \node[font=\footnotesize,align=center,text=figblue!90!black] at (6.55,3.13) {$\range T$\\[-1pt]\scriptsize $=(\nul T^*)^\perp$};

    \fill[black] (5.40,1.66) circle (1.8pt) node[right=3pt,font=\scriptsize] {$0_W$};
    \node[font=\scriptsize] at (6.75,1.66) {$\perp$};

    \draw[draw=figred!85!black,dashed,line width=0.6pt,fill=figred!5] (5.40,0.0) rectangle (7.70,1.25);
    \node[font=\footnotesize,text=figred!90!black] at (6.55,0.62) {$\nul T^*$};

    % Bijections between range T* and range T
    \draw[->,figblue,line width=0.8pt] (2.44,3.38) -- node[above=1pt,font=\scriptsize,text=figblue!90!black] {$T$ (bijection)} (5.26,3.38);
    \draw[->,figblue,line width=0.8pt] (5.26,2.86) -- node[below=1pt,font=\scriptsize,text=figblue!90!black] {$T^*$ (bijection)} (2.44,2.86);

    % Null spaces collapse cleanly to 0_W and 0_V with generous spacing
    \draw[->,figred!85!black,dashed,line width=0.65pt] (2.44,0.88) .. controls (3.65,1.12) and (4.45,1.42) .. (5.24,1.60);
    \node[font=\scriptsize,text=figred!90!black] at (3.85,2.04) {$T(\nul T)=\{0_W\}$};
    \draw[->,figred!85!black,dashed,line width=0.65pt] (5.26,0.46) .. controls (4.10,0.72) and (3.25,1.26) .. (2.46,1.56);
    \node[font=\scriptsize,text=figred!90!black] at (3.85,0.12) {$T^*(\nul T^*)=\{0_V\}$};
  \end{scope}

  % Right example: two side-by-side planes V and W cleanly separated from left schematic
  \begin{scope}[xshift=10.15cm,yshift=1.82cm]
    \node[font=\scriptsize,align=center,anchor=south] at (0,1.94) {$V=\RR^2$\\[-1pt]$\nul T \perp \range T^*$};
    \draw[figaxis,draw=black!35] (-1.22,0) -- (1.25,0);
    \draw[figaxis,draw=black!35] (0,-1.22) -- (0,1.28);
    \draw[very thick,figred,dashed,<->] (-1.12,0.56) -- (1.12,-0.56) node[below=2pt,font=\scriptsize,text=figred!90!black] {$\nul T$};
    \draw[very thick,figblue,<->] (-0.56,-1.12) -- (0.56,1.12) node[above right=1pt,font=\scriptsize,text=figblue!90!black] {$\range T^*$};
    % Right-angle mark in Q1/Q4 completely clear of axes
    \draw[black!70,thin] (0.11,0.22) -- (0.33,0.11) -- (0.22,-0.11);
  \end{scope}

  \begin{scope}[xshift=13.55cm,yshift=1.82cm]
    \node[font=\scriptsize,align=center,anchor=south] at (0,1.94) {$W=\RR^2$\\[-1pt]$\range T \perp \nul T^*$};
    \draw[very thick,figblue,<->] (-1.22,0) -- (1.25,0);
    \node[above=3pt,font=\scriptsize,text=figblue!90!black] at (0.76,0) {$\range T$};
    \draw[very thick,figred,dashed,<->] (0,-1.20) -- (0,1.25);
    \node[right=3pt,font=\scriptsize,text=figred!90!black] at (0,1.02) {$\nul T^*$};
    \draw[black!70,thin] (0.22,0) -- (0.22,0.22) -- (0,0.22);
  \end{scope}
\end{tikzpicture}
\figcap{\textbf{Left:} The four fundamental spaces of $T$ and $T^*$ (\hyperref[prop-94]{Prop.~94}): $V=\range T^*\oplus\nul T$ and $W=\range T\oplus\nul T^*$, where $T$ and $T^*$ restrict to bijections between $\range T^*$ and $\range T$ (blue) while collapsing $\nul T\to\{0_W\}$ and $\nul T^*\to\{0_V\}$ (red). \textbf{Right:} For $T(x,y)=(x+2y,0)$ on $\RR^2$ ($T^*(x,y)=(x,2x)$), $\range T$ (the $x$-axis) $\perp\nul T^*$ (the $y$-axis), and $\nul T=\spn(-2,1)\perp\range T^*=\spn(1,2)$.}
\end{fig}

\begin{proposition*}[LADR 7.7: Injectivity/surjectivity duality and ranks]\target{name-ladr-7-7}
  For $T\in\LL(V,W)$:
  \[
    T \text{ is injective} \iff T^* \text{ is surjective}; \qquad
    T \text{ is surjective} \iff T^* \text{ is injective}.
  \]
  Furthermore, $\dim\range T^* = \dim\range T$ and $\dim\nul T^* = \dim\nul T + \dim W - \dim V$.
\end{proposition*}
\begin{proof}
  \begin{steps}
    \ii $T$ is injective $\iff \nul T=\{0\} \iff (\nul T)^\perp=V$ \ilink{T5}{(T5)} $\iff \range T^*=V$ (Prop.~94(ii)) $\iff T^*$ is surjective.
    \ii $T$ is surjective $\iff \range T=W \iff (\range T)^\perp=\{0\}$ \ilink{T5}{(T5)} $\iff \nul T^*=\{0\}$ (Prop.~94(i)) $\iff T^*$ is injective.
    \ii By Prop.~94(ii), \ilink{T5}{(T5)}, and rank--nullity \ilink{T13}{(T13)}, $\dim\range T^* = \dim V - \dim\nul T = \dim\range T$.
    \ii By Prop.~94(i) and \ilink{T5}{(T5)}, $\dim\nul T^* = \dim W - \dim\range T = \dim W - (\dim V - \dim\nul T)$.
  \end{steps}
\end{proof}

\begin{proposition*}[LADR 7A: Eigenvalues of $T^*$]\target{name-ladr-7a-eigenvalues}
  For an operator $T\in\LL(V)$, $\lambda$ is an eigenvalue of $T$ if and only if $\overline{\lambda}$ is an eigenvalue of $T^*$.
\end{proposition*}
\begin{proof}
  By Prop.~94(i) and LADR~7.5, $\nul(T^*-\overline{\lambda}I)=\nul((T-\lambda I)^*)=\bigl(\range(T-\lambda I)\bigr)^\perp$. This is nonzero $\iff T-\lambda I$ is not surjective $\iff T-\lambda I$ is not injective (T13, since $T$ is an operator on a finite-dimensional space) $\iff \lambda$ is an eigenvalue of $T$.
\end{proof}

The eigenvectors of $T$ and $T^*$ need not coincide. For the shear matrix $B=\begin{psmallmatrix}1&1\\0&1\end{psmallmatrix}$, $Be_1=e_1$, while the $1$-eigenvectors of $B^{\mathsf T}=\begin{psmallmatrix}1&0\\1&1\end{psmallmatrix}$ are multiples of $e_2$.

\begin{shortcuts}
  \ii \textbf{Adjoint from a matrix:} In orthonormal bases of both domain and codomain, conjugate-transpose $\MM(T)$ (Prop.~93). Example: for $T(z_1,z_2)=(z_1+iz_2,3z_2)$ on $\CC^2$ with standard inner product, $T^*(z_1,z_2)=(z_1,-iz_1+3z_2)$.
  \ii \textbf{Adjoint in non-orthonormal bases:} Write $T^*(\text{basis vector})$ with unknown coefficients and enforce $\langle Tv,w\rangle=\langle v,T^*w\rangle$ (Ex.~72), or first apply Gram--Schmidt to orthonormalize (Ex.~73).
  \ii \textbf{Shifting by a scalar:} $(T-\lambda I)^* = T^* - \overline{\lambda}I$ (LADR~7.5). Example: $(T-iI)^*=T^*+iI$.
  \ii \textbf{Finding $\nul T^*$ directly:} $\nul T^* = (\range T)^\perp$ without computing $T^*$ (Prop.~94).
  \ii \textbf{Injective $\leftrightarrow$ surjective:} $T$ injective $\iff T^*$ surjective (LADR~7.7).
\end{shortcuts}

\begin{trapbox}
  \ii \textbf{Transpose requires orthonormal bases:} In the basis $1,x$, $\MM(D)^{\mathsf T}$ wrongly predicts $D^*(1)=x$ instead of $12x-6$.
  \ii \textbf{Conjugates over $\CC$:} $(\lambda T)^*=\overline{\lambda}T^*$, so $(iI)^* = -iI$.
  \ii \textbf{Products reverse order:} $(ST)^*=T^*S^*$, \emph{not} $S^*T^*$. For $S=\begin{psmallmatrix}0&1\\0&0\end{psmallmatrix}$ and $T=\begin{psmallmatrix}0&0\\1&0\end{psmallmatrix}$ on $\RR^2$, $(ST)^*=\operatorname{diag}(1,0)$ whereas $S^*T^*=\operatorname{diag}(0,1)$.
  \ii \textbf{Codomains swap:} $T\in\LL(V,W)$ gives $T^*\in\LL(W,V)$. Only for an operator ($V=W$) can $T$ and $T^*$ be added or compared.
\end{trapbox}

\subsection{Self-adjoint operators (Ch.~19)}

\begin{qbox}
  The matrix $A=\begin{psmallmatrix}2&1+i\\1-i&3\end{psmallmatrix}$ equals its own conjugate transpose ($A^*=A$). Must all of its eigenvalues be real? More generally, what numerical signature of $T$ detects whether $T=T^*$?
\end{qbox}

Among complex operators, those satisfying $T=T^*$ behave like the real numbers $\RR\subset\CC$ (where $z=\overline{z}$): their eigenvalues, quadratic values $\langle Tv,v\rangle$, and algebraic properties mirror real scalars.

\begin{definition*}[Def.~50 / LADR 7.10: Self-adjoint operator]\target{def-50}\target{name-self-adjoint}
  An operator $T\in\LL(V)$ is \textbf{self-adjoint} if $T=T^*$, i.e.\ $\langle Tv,w\rangle=\langle v,Tw\rangle$ for all $v,w\in V$. Over $\RR$ such operators (or matrices in an orthonormal basis) are called \textbf{symmetric}; over $\CC$ they are called \textbf{Hermitian} ($A=A^*$).
\end{definition*}

\begin{example*}[Ex.~74: Hermitian vs.\ complex symmetric]\target{ex-74}
  The first matrix below is Hermitian ($M_1^*=M_1$); the second equals its transpose $M_2^{\mathsf T}$ but is \emph{not} Hermitian ($M_2^*\neq M_2$), and its eigenvalues are $1\pm i\notin\RR$:
  \[
    M_1 = \begin{bmatrix}1&i\\-i&1\end{bmatrix} \text{ (Hermitian)}, \qquad
    M_2 = \begin{bmatrix}1&i\\i&1\end{bmatrix} \text{ (not Hermitian)}.
  \]
\end{example*}

\begin{theorem*}[Prop.~95 / LADR 7.12: Eigenvalues of self-adjoint operators are real]\target{prop-95}
  If $T\in\LL(V)$ is self-adjoint, then every eigenvalue of $T$ is real.
\end{theorem*}
\begin{ideabox}
  Evaluate $\langle Tv,v\rangle$ on a unit (or arbitrary nonzero) eigenvector $v$: pulling $\lambda$ from the first slot gives $\lambda\|v\|^2$, while moving $T=T^*$ to the second slot pulls out $\overline{\lambda}\|v\|^2$. Since $\|v\|^2\neq0$, $\lambda=\overline{\lambda}$.
\end{ideabox}
\begin{proof}
  Let $Tv=\lambda v$ with $v\neq0$. Using \ilink{T1}{(T1)} and $T=T^*$,
  \[
    \lambda\|v\|^2 = \langle \lambda v,v\rangle = \langle Tv,v\rangle = \langle v,T^*v\rangle = \langle v,Tv\rangle = \langle v,\lambda v\rangle = \overline{\lambda}\|v\|^2.
  \]
  Since $\|v\|^2>0$, dividing yields $\lambda=\overline{\lambda}$, so $\lambda\in\RR$.
\end{proof}

Notice the converse fails: $B=\begin{psmallmatrix}1&1\\0&1\end{psmallmatrix}$ has only the real eigenvalue $1$ yet is not self-adjoint.

\noindent\textbf{The question answered.} Yes: $A=A^*$ is Hermitian, so its eigenvalues are guaranteed to be real. Indeed by \ilink{T12}{(T12)}, $\tr A=5$ and $\det A=6-(1+i)(1-i)=4$, so $\lambda^2-5\lambda+4=(\lambda-1)(\lambda-4)$ gives eigenvalues $1$ and $4$.

\medskip\noindent\textbf{Polarization identity for operators.}\target{name-polarization} The ``diagonal'' quadratic values $Q(x)=\langle Tx,x\rangle$ determine the full sesquilinear form $\langle Tv,w\rangle$. Over $\CC$, for \emph{any} operator $T\in\LL(V)$ and all $v,w\in V$:
\[
  \langle Tv,w\rangle = \frac{1}{4}\bigl(Q(v+w)-Q(v-w)+iQ(v+iw)-iQ(v-iw)\bigr).
\]
\emph{Derivation.} Set $a=\langle Tv,w\rangle$ and $b=\langle Tw,v\rangle$. Expanding by \ilink{T1}{(T1)}:
\begin{align*}
  Q(v+w)-Q(v-w) &= (\langle Tv,v\rangle + a + b + \langle Tw,w\rangle) - (\langle Tv,v\rangle - a - b + \langle Tw,w\rangle) = 2a+2b,\\
  Q(v+iw)-Q(v-iw) &= 2\langle Tv,iw\rangle + 2\langle T(iw),v\rangle = -2ia+2ib.
\end{align*}
Multiplying the second identity by $i$ and adding gives $(2a+2b)+(2a-2b)=4a$. No hypothesis on $T$ is used. Over $\RR$, multiplication by $i$ is unavailable so only the first line applies; however, if $\FF=\RR$ and $T$ is self-adjoint, then $b=\langle Tw,v\rangle=\langle w,Tv\rangle=\langle Tv,w\rangle=a$, yielding $Q(v+w)-Q(v-w)=4\langle Tv,w\rangle$.

\begin{theorem*}[Prop.~96 / LADR 7.14: Over $\CC$, self-adjoint means $\langle Tv,v\rangle\in\RR$]\target{prop-96}
  Let $V$ be a complex inner product space and $T\in\LL(V)$. Then $T$ is self-adjoint if and only if $\langle Tv,v\rangle\in\RR$ for every $v\in V$.
\end{theorem*}
\begin{ideabox}
  One direction computes $\overline{\langle Tv,v\rangle}=\langle v,Tv\rangle$ directly when $T=T^*$. Conversely, if $Q(x)=\langle Tx,x\rangle$ is always real, swapping $v$ and $w$ in the polarization identity takes the complex conjugate of the expression for $\langle Tv,w\rangle$, proving $\langle Tw,v\rangle=\overline{\langle Tv,w\rangle}$ and thus $T=T^*$.
\end{ideabox}
\begin{proof}
  \begin{steps}
    \ii ($\Rightarrow$) If $T=T^*$, then $\overline{\langle Tv,v\rangle}=\langle v,Tv\rangle=\langle T^*v,v\rangle=\langle Tv,v\rangle$ by \ilink{T1}{(T1)}, so $\langle Tv,v\rangle\in\RR$.
    \ii ($\Leftarrow$) Suppose $Q(x)=\langle Tx,x\rangle\in\RR$ for all $x\in V$. Note $Q(\mu x)=|\mu|^2 Q(x)$, so $Q(-x)=Q(x)$ and $Q(\pm ix)=Q(x)$. Swapping $v$ and $w$ in polarization and rewriting $Q(w+iv)=Q(i(v-iw))=Q(v-iw)$ and $Q(w-iv)=Q(-i(v+iw))=Q(v+iw)$ gives
    \[
      \langle Tw,v\rangle = \frac{1}{4}\bigl(Q(v+w)-Q(v-w)+iQ(v-iw)-iQ(v+iw)\bigr) = \overline{\langle Tv,w\rangle},
    \]
    because every $Q(\cdot)$ is real. Hence $\langle Tv,w\rangle = \overline{\langle Tw,v\rangle} = \langle v,Tw\rangle$ for all $v,w$, so $T^*=T$ by \ilink{T4}{(T4)}.
  \end{steps}
\end{proof}

\begin{theorem*}[Prop.~97 / LADR 7.16, 7.13: Vanishing quadratic form]\target{prop-97}
  Suppose $\langle Tv,v\rangle=0$ for every $v\in V$. If either $V$ is complex or $T$ is self-adjoint, then $T=0$.
\end{theorem*}
\begin{proof}
  Here $Q\equiv 0$. Over $\CC$, the complex polarization identity gives $\langle Tv,w\rangle=0$ for all $v,w$. Over $\RR$ with $T$ self-adjoint, $4\langle Tv,w\rangle=Q(v+w)-Q(v-w)=0$ for all $v,w$. Taking $w=Tv$ in either case gives $\|Tv\|^2=0$, so $Tv=0$ for all $v\in V$, i.e.\ $T=0$.
\end{proof}

Over $\RR$, the self-adjoint hypothesis cannot be dropped: a $90^\circ$ rotation $R(x,y)=(-y,x)$ satisfies $\langle R(x,y),(x,y)\rangle = -yx+xy = 0$ for all $(x,y)\in\RR^2$, yet $R\neq0$ ($R^*=-R$).

\begin{shortcuts}
  \ii \textbf{Checking Hermitian:} Compare $\MM(T)$ with its conjugate transpose: diagonal entries must be real and $a_{ji}=\overline{a_{ij}}$ (Def.~50, Prop.~93).
  \ii \textbf{Products:} For self-adjoint $S,T$, $ST$ is self-adjoint $\iff ST=TS$ (since $(ST)^*=T^*S^*=TS$).
  \ii \textbf{Real combinations:} If $S,T$ are self-adjoint and $a,b\in\RR$, then $aS+bT$ is self-adjoint. Over $\CC$, $iT$ is skew-adjoint ($(iT)^*=-iT$).
  \ii \textbf{Quadratic test over $\CC$:} $T$ is Hermitian $\iff \langle Tv,v\rangle\in\RR\ \forall v$ (Prop.~96). For $A$ in the question, $\langle Av,v\rangle = 2|v_1|^2+3|v_2|^2+2\operatorname{Re}\bigl((1+i)v_2\overline{v_1}\bigr)\in\RR$.
\end{shortcuts}

\begin{trapbox}
  \ii \textbf{Real eigenvalues $\not\Rightarrow$ self-adjoint:} $\begin{psmallmatrix}1&1\\0&1\end{psmallmatrix}$ has only eigenvalue $1$ yet is not symmetric.
  \ii \textbf{Complex symmetric $\neq$ Hermitian:} $M_2$ in Ex.~74 satisfies $M_2^{\mathsf T}=M_2$, yet has eigenvalues $1\pm i$.
  \ii \textbf{$\langle Tv,v\rangle=0\ \forall v\ \not\Rightarrow T=0$ over $\RR$:} Counterexample: $90^\circ$ rotation $(x,y)\mapsto(-y,x)$.
\end{trapbox}

\subsection{Normal operators (Ch.~19)}

\begin{qbox}
  Compare the lengths $\|Ae_1\|$ and $\|A^*e_1\|$, as well as $\|Be_1\|$ and $\|B^*e_1\|$, for
  \[
    A=\begin{bmatrix}1&-1\\1&1\end{bmatrix}, \qquad B=\begin{bmatrix}1&1\\0&1\end{bmatrix}.
  \]
  Which operators satisfy $\|Tv\|=\|T^*v\|$ for \emph{every} vector $v\in V$?
\end{qbox}

Self-adjointness ($T=T^*$) is too strict to include rotations or complex operators with non-real eigenvalues, even though rotations have an orthonormal eigenbasis over $\CC$. The exact condition giving an orthonormal eigenbasis over $\CC$ is commutativity with the adjoint.

\begin{definition*}[Def.~51 / LADR 7A: Normal operator]\target{def-51}\target{name-normal}
  An operator $T\in\LL(V)$ is \textbf{normal} if it commutes with its adjoint:
  \[
    TT^* = T^*T.
  \]
\end{definition*}

Every self-adjoint operator is normal ($TT^*=T^2=T^*T$). The matrix $A$ in the question is normal without being self-adjoint:
\[
  AA^* = \begin{bmatrix}1&-1\\1&1\end{bmatrix}\begin{bmatrix}1&1\\-1&1\end{bmatrix} = \begin{bmatrix}2&0\\0&2\end{bmatrix} = A^*A.
\]
By contrast, $B$ is not normal: $(BB^*)_{11}=2\neq 1=(B^*B)_{11}$.

\begin{proposition*}[LADR 7.23: Real and imaginary parts of a complex operator]\target{name-ladr-7-23}
  Over $\CC$, every operator $T\in\LL(V)$ splits uniquely as $T=R+iS$ with $R,S$ self-adjoint, namely
  \[
    R = \frac{T+T^*}{2}, \qquad S = \frac{T-T^*}{2i}.
  \]
  Moreover, $T$ is normal if and only if $R$ and $S$ commute: $RS=SR$.
\end{proposition*}
\begin{ideabox}
  Just as a complex number $z=x+iy$ decomposes uniquely into its real part $(z+\overline{z})/2$ and imaginary part $(z-\overline{z})/(2i)$, any operator $T$ decomposes into self-adjoint ``real'' and ``imaginary'' parts $R$ and $S$. Taking adjoints gives $T^*=R-iS$, so comparing $TT^*=(R+iS)(R-iS)=R^2+S^2-i(RS-SR)$ with $T^*T=R^2+S^2+i(RS-SR)$ immediately shows $TT^*=T^*T \iff RS=SR$.
\end{ideabox}
\begin{proof}
  \begin{steps}
    \ii Clearly $R+iS=T$. By LADR~7.5, $R^*=\frac{T^*+T}{2}=R$ and $S^*=\frac{T^*-T}{-2i}=S$, so both are self-adjoint. If $T=R_1+iS_1$ with $R_1,S_1$ self-adjoint, then $T^*=R_1-iS_1$, forcing $R_1=(T+T^*)/2=R$ and $S_1=(T-T^*)/(2i)=S$.
    \ii Expanding using $R=R^*,S=S^*$ gives
    \[
      T^*T - TT^* = (R-iS)(R+iS) - (R+iS)(R-iS) = 2i(RS-SR),
    \]
    so $RS-SR = \frac{T^*T-TT^*}{2i}$, which vanishes if and only if $T^*T=TT^*$.
  \end{steps}
\end{proof}

Warning: $R$ and $S$ are \emph{not} obtained by taking the entrywise real and imaginary parts of $\MM(T)$. For the real matrix $T=\begin{psmallmatrix}0&-1\\1&0\end{psmallmatrix}$, $T+T^*=0$ so $R=0$, whereas $S=\begin{psmallmatrix}0&i\\-i&0\end{psmallmatrix}$ has purely imaginary entries.

\begin{theorem*}[Prop.~98 / LADR 7.20: Normality $\iff \|Tv\|=\|T^*v\|$]\target{prop-98}
  An operator $T\in\LL(V)$ is normal if and only if $\|Tv\|=\|T^*v\|$ for every $v\in V$.
\end{theorem*}
\begin{ideabox}
  Measure the failure of normality by the operator $N=T^*T-TT^*$. This defect $N$ is always self-adjoint, so $N=0$ is tested by the diagonal values $\langle Nv,v\rangle$ alone (Prop.~97), and moving $T^*$ or $T$ across reveals $\langle Nv,v\rangle = \|Tv\|^2-\|T^*v\|^2$.
\end{ideabox}
\begin{proof}
  \begin{steps}
    \ii Let $N=T^*T-TT^*$. By LADR~7.5, $N^*=(T^*T)^*-(TT^*)^* = T^*T-TT^*=N$, so $N$ is self-adjoint.
    \ii For every $v\in V$,
    \[
      \langle Nv,v\rangle = \langle T^*Tv,v\rangle - \langle TT^*v,v\rangle = \langle Tv,Tv\rangle - \langle T^*v,T^*v\rangle = \|Tv\|^2 - \|T^*v\|^2.
    \]
    \ii Thus $T$ is normal $\iff N=0 \iff \langle Nv,v\rangle=0\ \forall v\in V$ (Prop.~97) $\iff \|Tv\|=\|T^*v\|\ \forall v\in V$.
  \end{steps}
\end{proof}

\begin{fig}
\begin{tikzpicture}[scale=1.0,>=Stealth]
  % Left panel title aligned with right panel title
  \node[font=\footnotesize\bfseries] at (2.35,5.88) {(a) Directional stretch $r(\theta)$};

  \begin{axis}[
    at={(0.15cm,0.80cm)},
    anchor=south west,
    width=5.8cm,height=5.3cm,
    axis lines=middle,
    axis line style={-{Stealth[length=4pt,width=2.6pt]},black!75},
    xmin=-2.15,xmax=2.25,ymin=-2.05,ymax=2.25,
    xtick=\empty,ytick=\empty,
    xlabel={$x$},ylabel={$y$},
    label style={font=\scriptsize},
    axis equal image,
    clip=false
  ]
    \addplot[domain=0:360,samples=160,smooth,figslate,thick,densely dotted]
      ({sqrt(2)*cos(x)},{sqrt(2)*sin(x)});
    \addplot[domain=0:360,samples=200,smooth,figblue,very thick]
      ({sqrt((cos(x)+sin(x))^2 + sin(x)^2)*cos(x)},{sqrt((cos(x)+sin(x))^2 + sin(x)^2)*sin(x)});
    \addplot[domain=0:360,samples=200,smooth,figred,very thick,dashed]
      ({sqrt(cos(x)^2 + (cos(x)+sin(x))^2)*cos(x)},{sqrt(cos(x)^2 + (cos(x)+sin(x))^2)*sin(x)});

    % Direct curve callouts pointing cleanly to each curve's outermost lobe
    \node[font=\scriptsize,text=figslate] at (axis cs:0,-2.35) {$\|Av\|=\|A^*v\|=\sqrt{2}$ (dotted)};
    \node[font=\scriptsize,text=figblue!90!black] (bvlbl) at (axis cs:1.58,1.96) {$\|Bv\|$ (solid)};
    \draw[->,figblue!85!black,thin] (bvlbl.south) -- (axis cs:0.86,1.46);
    \node[font=\scriptsize,text=figred!90!black,anchor=east] (bslbl) at (axis cs:-1.18,-1.64) {$\|B^*v\|$ (dashed)};
    \draw[->,figred!85!black,thin] ([xshift=-12pt]bslbl.north east) -- (axis cs:-1.49,-0.78);
  \end{axis}

  % Right: Hierarchy of operator classes
  \begin{scope}[xshift=6.55cm,yshift=0.0cm]
    \node[font=\footnotesize\bfseries] at (4.00,5.88) {(b) Hierarchy of operator classes};
    \tikzset{
      opbox/.style={draw=black,line width=0.55pt,inner sep=4.2pt,align=center,font=\scriptsize},
      inc/.style={->,line width=0.7pt,shorten <=2pt,shorten >=2pt}
    }
    % Top: All operators L(V)
    \node[opbox,fill=black!3] (all) at (4.00,5.08)
      {\textbf{All operators $\LL(V)$}\\[1pt]e.g.\ shear $B=\begin{psmallmatrix}1&1\\0&1\end{psmallmatrix}$ ($BB^*\ne B^*B$)};

    % Level 2: Normal operators
    \node[opbox,draw=figblue!85!black,fill=figblue!5,text=figblue!90!black] (norm) at (4.00,3.50)
      {\textbf{Normal} ($TT^*=T^*T \iff \|Tv\|=\|T^*v\|$)\\[1pt]e.g.\ $A=\begin{psmallmatrix}1&-1\\1&1\end{psmallmatrix}$};

    % Level 3 Left: Self-adjoint
    \node[opbox,draw=figgreen!85!black,fill=figgreen!5,text=figgreen!90!black] (sa) at (1.80,1.95)
      {\textbf{Self-adjoint} ($T^*=T$)\\[1pt]real spectrum; $D=\begin{psmallmatrix}2&0\\0&-1\end{psmallmatrix}$};

    % Level 3 Right: Unitary / Isometries
    \node[opbox,draw=figamber!90!black,fill=figamber!5,text=figamber!90!black] (unit) at (6.20,1.95)
      {\textbf{Unitary / Isometry} ($T^*T=I$)\\[1pt]$\|Tv\|=\|v\|$; rotation $R_{\pi/2}$};

    % Level 4 Left: Positive (centered directly below Self-adjoint)
    \node[opbox,draw=figred!85!black,fill=figred!6,text=figred!90!black] (pos) at (1.55,0.38)
      {\textbf{Positive} ($\langle Tv,v\rangle\ge 0$)\\[1pt]$\lambda_j\ge 0$; $P=\begin{psmallmatrix}2&1\\1&2\end{psmallmatrix}$};

    % Level 4 Right: Intersection Self-adjoint \cap Unitary
    \node[opbox,draw=figslate,dashed,fill=black!3] (inter) at (5.85,0.38)
      {\textbf{Self-adjoint $\cap$ Unitary}\\[1pt]$T^*=T=T^{-1}$: reflections $H_{\text{refl}},\,\pm I$};

    % Clean inclusion arrows connecting node borders directly (zero floating gaps!)
    \draw[inc] (norm.north) -- (all.south);
    \draw[inc,figgreen!85!black] (sa.north) -- ($(norm.south)+(-1.05,0)$);
    \draw[inc,figamber!90!black] (unit.north) -- ($(norm.south)+(1.05,0)$);
    \draw[inc,figred!85!black] (pos.north) -- ($(sa.south)+(-0.25,0)$);
    \draw[inc,figslate,dashed] ($(inter.north)+(-1.10,0)$) -- ($(sa.south)+(0.80,0)$);
    \draw[inc,figslate,dashed] ($(inter.north)+(0.35,0)$) -- (unit.south);
  \end{scope}
\end{tikzpicture}
\figcap{\textbf{Left:} Polar plot of directional stretch $r(\theta)=\|Tv\|$ and $\|T^*v\|$ on unit vectors $v=(\cos\theta,\sin\theta)$ (\hyperref[prop-98]{Prop.~98}). Normal $A=\begin{psmallmatrix}1&-1\\1&1\end{psmallmatrix}$ produces identical circles of radius $\sqrt{2}$ (dotted); non-normal $B=\begin{psmallmatrix}1&1\\0&1\end{psmallmatrix}$ gives distinct curves for $\|Bv\|$ (blue solid) and $\|B^*v\|$ (red dashed). \textbf{Right:} Hierarchy of operator classes inside $\LL(V)$ with $2\times2$ examples.}
\end{fig}

\noindent\textbf{The question answered.} $\|Ae_1\|=\|(1,1)\|=\sqrt{2}=\|A^*e_1\|$, whereas $\|Be_1\|=1\neq \sqrt{2}=\|B^*e_1\|$. By Prop.~98, $\|Tv\|=\|T^*v\|$ holds for all $v$ precisely when $T$ is normal.

\begin{theorem*}[Prop.~99 / LADR 7.21: Null spaces, ranges, and eigenvectors of normal operators]\target{prop-99}
  Let $T\in\LL(V)$ be normal. Then:
  \begin{enumerate}
    \ii[(a)] $\nul T^* = \nul T$ and $\range T^* = \range T$, so $(\nul T)^\perp = \range T$ and $V = \nul T\oplus \range T$;
    \ii[(b)] $T-\lambda I$ is normal for every $\lambda\in\FF$;
    \ii[(c)] $v$ is an eigenvector of $T$ with eigenvalue $\lambda$ $\iff$ $v$ is an eigenvector of $T^*$ with eigenvalue $\overline{\lambda}$.
  \end{enumerate}
\end{theorem*}
\begin{proof}
  \begin{steps}
    \ii $\nul T^*=\nul T$: by Prop.~98, $T^*v=0 \iff \|T^*v\|=0 \iff \|Tv\|=0 \iff Tv=0$.
    \ii $\range T^*=\range T$: Prop.~94(ii) and Step~1 give $\range T^* = (\nul T)^\perp = (\nul T^*)^\perp$. By Prop.~94(i) and \ilink{T5}{(T5)}, $(\nul T^*)^\perp = \bigl((\range T)^\perp\bigr)^\perp = \range T$.
    \ii Direct sum: $V = \nul T\oplus(\nul T)^\perp = \nul T\oplus\range T$.
    \ii Shifting preserves normality: using LADR~7.5 and $TT^*=T^*T$,
    \[
      (T-\lambda I)(T-\lambda I)^* = (T-\lambda I)(T^*-\overline{\lambda}I) = TT^*-\overline{\lambda}T-\lambda T^*+|\lambda|^2I = (T^*-\overline{\lambda}I)(T-\lambda I).
    \]
    \ii Eigenvectors: applying Step~1 to the normal operator $T-\lambda I$ gives $\nul(T-\lambda I)=\nul(T^*-\overline{\lambda}I)$.
  \end{steps}
\end{proof}

Without normality, $V=\nul T\oplus\range T$ can fail completely: for $N=\begin{psmallmatrix}0&1\\0&0\end{psmallmatrix}$, $\nul N = \range N = \spn(e_1)$.

\begin{proposition*}[LADR 7.22: Orthogonality of eigenspaces for normal operators]\target{name-ladr-7-22}
  If $T\in\LL(V)$ is normal and $Tv=\lambda v$, $Tw=\mu w$ with $\lambda\neq\mu$, then $\langle v,w\rangle=0$.
\end{proposition*}
\begin{proof}
  By Prop.~99(c), $T^*w=\overline{\mu}w$. Moving $T$ across (Def.~49, T1) gives
  \[
    \lambda\langle v,w\rangle = \langle Tv,w\rangle = \langle v,T^*w\rangle = \langle v,\overline{\mu}w\rangle = \mu\langle v,w\rangle,
  \]
  so $(\lambda-\mu)\langle v,w\rangle=0$. Since $\lambda\neq\mu$, $\langle v,w\rangle=0$.
\end{proof}

Because self-adjoint operators are normal, eigenvectors of a symmetric/Hermitian operator for distinct eigenvalues are orthogonal. Non-normal operators lack this: $B=\begin{psmallmatrix}1&1\\0&2\end{psmallmatrix}$ has eigenvectors $(1,0)$ and $(1,1)$ which are not orthogonal.

\begin{shortcuts}
  \ii \textbf{Testing normality:} Compute $AA^*$ and $A^*A$ in any orthonormal basis and compare (Def.~51, Prop.~93).
  \ii \textbf{Eigenvectors of $T^*$:} For normal $T$, $T$ and $T^*$ share the exact same eigenvectors, with eigenvalues conjugated (Prop.~99). Example: for $A=\begin{psmallmatrix}1&-1\\1&1\end{psmallmatrix}$, $A(1,-i)=(1+i)(1,-i)$ forces $A^*(1,-i)=(1-i)(1,-i)$.
  \ii \textbf{Automatic normality:} Self-adjoint ($T^*=T$), skew-adjoint ($T^*=-T$), and isometric ($T^*T=I$, hence $T^*=T^{-1}$ by T13 so $TT^*=I$) operators are always normal.
  \ii \textbf{Null space and range of normal operators:} For normal $T$, $\nul T^*=\nul T$ and $(\nul T)^\perp=\range T$ (Prop.~99).
\end{shortcuts}

\begin{trapbox}
  \ii \textbf{Normal $\neq$ self-adjoint:} $A=\begin{psmallmatrix}1&-1\\1&1\end{psmallmatrix}$ is normal yet has non-real eigenvalues $1\pm i$ and $A^*\neq A$.
  \ii \textbf{Normal $+$ normal need not be normal:} adding a symmetric and a rotation matrix can be non-normal, since the sum $N$ fails $\|Ne_1\|=\|N^*e_1\|$ ($2\neq0$):
  \[
    \begin{bmatrix}0&1\\1&0\end{bmatrix}+\begin{bmatrix}0&-1\\1&0\end{bmatrix}=\begin{bmatrix}0&0\\2&0\end{bmatrix} = N.
  \]
  \ii \textbf{$\nul T^*=\nul T$ requires normality:} For $N=\begin{psmallmatrix}0&1\\0&0\end{psmallmatrix}$, $\nul N=\spn(e_1)$ whereas $\nul N^{\mathsf T}=\spn(e_2)$.
\end{trapbox}

\subsection{The real spectral theorem (Ch.~20)}

\begin{qbox}
  The matrix $A=\begin{psmallmatrix}3&1\\1&3\end{psmallmatrix}$ is real symmetric. Does $\RR^2$ have an \emph{orthonormal} basis of eigenvectors of $A$? Geometrically, how does $A$ transform the unit circle?
\end{qbox}

Throughout this subsection and the next, assume $V\neq\{0\}$.

\begin{definition*}[Orthogonally diagonalizable]
  An operator $T\in\LL(V)$ is \textbf{orthogonally diagonalizable} if $V$ admits an orthonormal basis consisting of eigenvectors of $T$.
\end{definition*}

Over $\CC$, every operator has an eigenvalue by the Fundamental Theorem of Algebra \ilink{T7}{(T7)}. To see why a real self-adjoint operator $T$ must have a \emph{real} eigenvalue, consider factoring $m_T$ via \ilink{T9}{(T9)}: any factor $x^2+bx+c$ with $b^2<4c$ produces an operator $T^2+bT+cI$ that is always injective (and hence invertible) when $T$ is self-adjoint, ruling out irreducible quadratic factors without real roots.

\begin{proposition*}[LADR 7.26: Invertibility of positive quadratic expressions]\target{name-ladr-7-26}
  Let $T\in\LL(V)$ be self-adjoint and $b,c\in\RR$ satisfy $b^2<4c$. Then $T^2+bT+cI$ is invertible.
\end{proposition*}
\begin{ideabox}
  For a real variable $t$, $t^2+bt+c>0$ whenever $b^2<4c$ by completing the square. Evaluate $\langle (T^2+bT+cI)v,v\rangle$ on $v\neq0$: self-adjointness turns $\langle T^2v,v\rangle$ into $\|Tv\|^2$, and Cauchy--Schwarz bounds $\langle Tv,v\rangle$ so that completing the square forces $\langle (T^2+bT+cI)v,v\rangle > 0$. Because the inner product is strictly positive, $(T^2+bT+cI)v$ cannot be $0$.
\end{ideabox}
\begin{proof}
  \begin{steps}
    \ii Fix $v\neq0$. Because $T=T^*$, $\langle T^2v,v\rangle = \langle Tv,Tv\rangle = \|Tv\|^2$, so
    \[
      \langle(T^2+bT+cI)v,v\rangle = \|Tv\|^2 + b\langle Tv,v\rangle + c\|v\|^2.
    \]
    \ii By Cauchy--Schwarz \ilink{T2}{(T2)}, $b\langle Tv,v\rangle \ge -|b|\,\|Tv\|\,\|v\|$. Completing the square yields
    \[
      \langle(T^2+bT+cI)v,v\rangle \ge \left(\|Tv\| - \frac{|b|}{2}\|v\|\right)^2 + \left(c - \frac{b^2}{4}\right)\|v\|^2 > 0,
    \]
    since the square is $\ge0$, $c-b^2/4>0$, and $\|v\|^2>0$ (positive definiteness, T1).
    \ii Hence $(T^2+bT+cI)v\neq0$ for every $v\neq0$. Thus $T^2+bT+cI$ is injective, hence invertible \ilink{T13}{(T13)}.
  \end{steps}
\end{proof}

\begin{theorem*}[Prop.~100 / LADR 7.27: Existence of eigenvalues for self-adjoint operators]\target{prop-100}
  Every self-adjoint operator $T\in\LL(V)$ on $V\neq\{0\}$ has an eigenvalue (over $\RR$ as well as over $\CC$).
\end{theorem*}
\begin{proof}
  Over $\CC$, this is \ilink{T7}{(T7)}. Suppose $\FF=\RR$ and suppose for contradiction that $m_T$ has no real root. By \ilink{T9}{(T9)}, $m_T=p_1\cdots p_k$ with each $p_j(x)=x^2+b_jx+c_j$ satisfying $b_j^2<4c_j$. By LADR~7.26, every $p_j(T)$ is invertible, so $m_T(T)=p_1(T)\cdots p_k(T)$ is invertible as a product of invertible operators. This contradicts $m_T(T)=0$ on $V\neq\{0\}$. Thus $m_T$ has a real root $\lambda$, which is an eigenvalue of $T$ by \ilink{T8}{(T8)}.
\end{proof}

\begin{theorem*}[Prop.~101 / LADR 7B.19: Orthogonal complements of invariant subspaces]\target{prop-101}
  Let $T\in\LL(V)$ be self-adjoint and let $U\subseteq V$ be a $T$-invariant subspace. Then $U^\perp$ is also $T$-invariant, and both restrictions $T|_U\in\LL(U)$ and $T|_{U^\perp}\in\LL(U^\perp)$ are self-adjoint (under the restricted inner product).
\end{theorem*}
\begin{proof}
  \begin{steps}
    \ii $U^\perp$ is $T$-invariant: if $v\in U^\perp$ and $u\in U$, then $Tu\in U$ (invariance of $U$), so $\langle Tv,u\rangle=\langle v,Tu\rangle=0$. Because $u\in U$ is arbitrary, $Tv\in U^\perp$.
    \ii $T|_U$ is self-adjoint: for $u,u'\in U$, $\langle T|_U u,u'\rangle=\langle Tu,u'\rangle=\langle u,Tu'\rangle=\langle u,T|_U u'\rangle$.
    \ii $T|_{U^\perp}$ is self-adjoint: by Step~1, $T|_{U^\perp}$ maps $U^\perp$ into $U^\perp$, and the same inner product identity holds for all $x,x'\in U^\perp$.
  \end{steps}
\end{proof}

Because $T$ is self-adjoint, \emph{every} $T$-invariant subspace has an orthogonal $T$-invariant complement $U^\perp$. Non-normal operators can fail this completely: $\begin{psmallmatrix}1&1\\0&1\end{psmallmatrix}$ has $\spn(e_1)$ as its only $1$-dimensional invariant subspace, so no invariant complement exists at all.

\begin{theorem*}[Prop.~102 / LADR 7.29: The real spectral theorem]\target{prop-102}
  Let $T\in\LL(V)$.
  \begin{enumerate}
    \ii[(a)] Every self-adjoint operator on $V$ is orthogonally diagonalizable (over $\RR$ or $\CC$).
    \ii[(b)] If $\FF=\RR$, conversely every orthogonally diagonalizable operator is self-adjoint.
  \end{enumerate}
  Thus over $\RR$: $T$ is self-adjoint $\iff$ $T$ is orthogonally diagonalizable.
\end{theorem*}
\begin{ideabox}
  Prove (a) by splitting off one eigendirection at a time. By Prop.~100, $T$ has an eigenvector $v_1$ (rescale to $\|v_1\|=1$). Its line $U=\spn(v_1)$ is $T$-invariant, so by Prop.~101 its orthogonal complement $U^\perp$ of dimension $n-1$ is also $T$-invariant and $T|_{U^\perp}$ is self-adjoint. By induction $U^\perp$ has an orthonormal eigenbasis $v_2,\dots,v_n$, and adjoining $v_1$ gives one for $V$. (For (b), test $\langle Tv_i,v_j\rangle=\langle v_i,Tv_j\rangle$ in the orthonormal eigenbasis: real eigenvalues factor out of the second slot unconjugated.)
\end{ideabox}
\begin{proof}
  \begin{steps}
    \ii \emph{Proof of (a) by induction on $n=\dim V$.} For $n=1$, any unit vector is an orthonormal eigenbasis. Assume $n>1$. By Prop.~100, $T$ has an eigenvector; scale it to a unit vector $v_1$. Set $U=\spn(v_1)$, which is $T$-invariant. By Prop.~101 and \ilink{T5}{(T5)}, $U^\perp$ is $T$-invariant of dimension $n-1$ and $T|_{U^\perp}$ is self-adjoint. By induction, $U^\perp$ has an orthonormal eigenbasis $v_2,\dots,v_n$ for $T|_{U^\perp}$ (hence eigenvectors of $T$). Since $v_1\perp U^\perp$, $v_1,v_2,\dots,v_n$ is an orthonormal list of $n$ vectors, hence a basis \ilink{T3}{(T3)} of eigenvectors of $T$.
    \ii \emph{Proof of (b) over $\RR$.} Let $v_1,\dots,v_n$ be an orthonormal basis with $Tv_i=\lambda_i v_i$ ($\lambda_i\in\RR$). For all $i,j$,
    \[
      \langle Tv_i,v_j\rangle = \lambda_i\delta_{ij} = \lambda_j\delta_{ij} = \langle v_i,\lambda_j v_j\rangle = \langle v_i,Tv_j\rangle,
    \]
    because $\lambda_i\delta_{ij}=\lambda_j\delta_{ij}=0$ when $i\neq j$, and $\lambda_j\in\RR$. By linearity in each slot, $\langle Tv,w\rangle=\langle v,Tw\rangle$ for all $v,w\in V$, so $T=T^*$.
  \end{steps}
\end{proof}

\begin{proposition*}[LADR 7.27: Minimal polynomial of a self-adjoint operator]\target{name-ladr-7-27}
  If $T\in\LL(V)$ is self-adjoint, then $m_T = (x-\lambda_1)\cdots(x-\lambda_m)$, where $\lambda_1,\dots,\lambda_m\in\RR$ are the distinct eigenvalues of $T$.
\end{proposition*}
\begin{proof}
  By Prop.~102(a), $T$ is diagonalizable, so $m_T$ is a product of distinct linear factors \ilink{T8}{(T8)} whose roots are the eigenvalues of $T$, which are real by Prop.~95.
\end{proof}

\noindent\textbf{Matrix form ($A=QDQ^{\mathsf T}$).} A real square matrix $Q$ with orthonormal columns satisfies $Q^{\mathsf T}Q=I$ (the $(i,j)$-entry is $q_i\cdot q_j=\delta_{ij}$); such a matrix is called \textbf{orthogonal}, and $Q^{-1}=Q^{\mathsf T}$. If $A$ is real symmetric and $Q=[u_1\;\cdots\;u_n]$ holds its orthonormal eigenvectors from Prop.~102 with $D=\operatorname{diag}(\lambda_1,\dots,\lambda_n)$, then $AQ=QD$ column-by-column, so
\[
  A = QDQ^{\mathsf T}, \qquad A^m = QD^m Q^{\mathsf T}\quad (m\ge1).
\]

\noindent\textbf{Recipe for orthogonal diagonalization of a real symmetric matrix $A$.}
\begin{steps}
  \ii Find all eigenvalues (guaranteed real) and a basis for each eigenspace.
  \ii Apply Gram--Schmidt \ilink{T3}{(T3)} \emph{inside each eigenspace} of dimension $\ge2$ (vectors from distinct eigenspaces are automatically orthogonal by LADR~7.22).
  \ii Place the resulting unit vectors as columns of $Q$ and matching eigenvalues on $D$; then $A=QDQ^{\mathsf T}$.
\end{steps}

\begin{example*}[Worked example: repeated eigenvalue in $3\times3$]
  Consider the symmetric matrix
  \[
    A = \begin{bmatrix}2&1&1\\1&2&1\\1&1&2\end{bmatrix}.
  \]
  Observe $A-I=J$ (all $1$s), which has rank $1$, so $\dim E(1,A)=3-1=2$ \ilink{T13}{(T13)}. Meanwhile $A(1,1,1)=4(1,1,1)$, so $\dim E(4,A)\ge1$. Because eigenspaces form a direct sum, $\dim E(4,A)=1$ and there are no other eigenvalues ($4,1,1$). By LADR~7.27, $m_A=(x-4)(x-1)$ without repeated roots (check $(A-4I)(A-I)=(J-3I)J=0$ since $J^2=3J$).

  A basis for the $1$-eigenspace $x+y+z=0$ is $w_1=(1,-1,0)$, $w_2=(1,0,-1)$, whose dot product is $1\neq0$. Applying Gram--Schmidt gives $e_2=\frac{1}{\sqrt{2}}(1,-1,0)$ and
  \[
    v_3 = (1,0,-1) - \langle(1,0,-1),e_2\rangle e_2 = (1,0,-1)-\frac12(1,-1,0)=\frac12(1,1,-2), \qquad e_3=\frac{1}{\sqrt{6}}(1,1,-2).
  \]
  With $e_1=\frac{1}{\sqrt{3}}(1,1,1)\perp E(1,A)$, we obtain $A=QDQ^{\mathsf T}$ where
  \[
    Q = \begin{bmatrix}1/\sqrt{3}&1/\sqrt{2}&1/\sqrt{6}\\1/\sqrt{3}&-1/\sqrt{2}&1/\sqrt{6}\\1/\sqrt{3}&0&-2/\sqrt{6}\end{bmatrix}, \qquad
    D = \begin{bmatrix}4&0&0\\0&1&0\\0&0&1\end{bmatrix}.
  \]
\end{example*}

\begin{fig}
\begin{tikzpicture}[scale=1.0,>=Stealth]
  \begin{scope}[xshift=0cm,yshift=0cm,scale=0.55]
    % Proper axes with arrowheads and legible labels
    \draw[figaxis,draw=black!50] (-3.8,0) -- (4.2,0) node[right=2pt,font=\scriptsize,text=black] {$x$};
    \draw[figaxis,draw=black!50] (0,-3.6) -- (0,3.8) node[above=1pt,font=\scriptsize,text=black] {$y$};

    % Image ellipse with light tint
    \fill[figblue,fill opacity=0.05] (0,0) [rotate=45] ellipse (4 and 2);
    \draw[figblue,very thick,rotate=45] (0,0) ellipse (4 and 2);

    % Unit circle (dashed) with label cleanly inside open Q3 interior (clear of ellipse!)
    \draw[black!65,thin,dashed] (0,0) circle (1);
    \node[font=\scriptsize,text=black!75,anchor=north east] at (-0.82,-0.82) {$\|v\|=1$};

    % Right-angle mark between u1 and u2
    \draw[black!65,thin,rotate=45] (0.42,0) -- (0.42,-0.42) -- (0,-0.42);

    % Stretched eigenvectors 4u1 and 2u2
    \draw[->,very thick,figblue,shorten >=1pt] (0,0) -- ({4/sqrt(2)},{4/sqrt(2)})
      node[above right=1pt,font=\scriptsize,text=figblue!90!black] {$4u_1$};
    \draw[->,very thick,figred,shorten >=1pt] (0,0) -- ({2/sqrt(2)},{-2/sqrt(2)})
      node[below right=1pt,font=\scriptsize,text=figred!90!black] {$2u_2$};

    % Unit eigenvectors u1 and u2 cleanly labeled in open Q1/Q4 space (clear of y-axis & unit circle!)
    \draw[->,thick,black] (0,0) -- ({1/sqrt(2)},{1/sqrt(2)});
    \node[font=\scriptsize,anchor=west] (u1lbl) at (1.28,0.52) {$u_1$};
    \draw[->,black!60,thin] (u1lbl.west) -- (0.78,0.68);
    \draw[->,thick,black] (0,0) -- ({1/sqrt(2)},{-1/sqrt(2)});
    \node[font=\scriptsize,anchor=west] (u2lbl) at (1.28,-0.52) {$u_2$};
    \draw[->,black!60,thin] (u2lbl.west) -- (0.78,-0.68);

    \node[font=\footnotesize,anchor=north] at (0,-3.95) {(a) Unit circle $\mapsto$ ellipse ($u_1\perp u_2$)};
  \end{scope}
\end{tikzpicture}
\hspace{0.9cm}
\tdplotsetmaincoords{68}{118}
\begin{tikzpicture}[tdplot_main_coords,scale=1.10,>=Stealth]
  % Exact orthonormal frame for E(1,A) = span(e2,e3) (z=0) and E(4,A) = span(e1) (z-axis)
  \draw[figred,thick,dashed] (0,0,-1.05) -- (0,0,0);

  \fill[figplane] (-1.55,-1.55,0) -- (2.05,-1.55,0) -- (2.05,2.05,0) -- (-1.55,2.05,0) -- cycle;
  \draw[figblue!70!black,thin] (-1.55,-1.55,0) -- (2.05,-1.55,0) -- (2.05,2.05,0) -- (-1.55,2.05,0) -- cycle;

  % Crisp right-angle marks showing e2 \perp e3 and e1 \perp E(1,A)
  \draw[figblue!70!black,thin] (0.28,0,0) -- (0.28,0.28,0) -- (0,0.28,0);
  \draw[black!70,thin] (0.28,0,0) -- (0.28,0,0.28) -- (0,0,0.28);
  \draw[black!70,thin] (0,0.28,0) -- (0,0.28,0.28) -- (0,0,0.28);

  % Orthonormal basis e2, e3 of E(1,A)
  \draw[->,very thick,figblue] (0,0,0) -- (1.45,0,0);
  \node[font=\scriptsize,text=figblue!90!black,above left=2pt] at (1.25,0,0) {$e_2$};
  \draw[->,very thick,figblue] (0,0,0) -- (0,1.45,0);
  \node[font=\scriptsize,text=figblue!90!black,above right=2pt] at (0,1.25,0) {$e_3$};

  % Plane label cleanly inside open front interior (no halo!)
  \node[font=\scriptsize,text=figblue!90!black] at (1.35,1.35,0) {$E(1,A)$};

  % Normal eigenvector e1 in E(4,A)
  \draw[->,very thick,figred] (0,0,0) -- (0,0,1.85);
  \node[font=\scriptsize,text=figred!90!black,right=5pt] at (0,0,1.58) {$e_1\in E(4,A)$};

  \fill[black] (0,0,0) circle (1.4pt);

  \begin{scope}[tdplot_screen_coords]
    \node[font=\footnotesize,anchor=north] at (0,-1.55) {(b) Orthogonal eigenspaces $E(4,A)\perp E(1,A)$};
  \end{scope}
\end{tikzpicture}
\figcap{\textbf{Left:} For $A=\begin{psmallmatrix}3&1\\1&3\end{psmallmatrix}$, the unit circle maps to an ellipse whose perpendicular semi-axes lie along $u_1=\frac{1}{\sqrt{2}}(1,1)$ (stretch $4$, blue) and $u_2=\frac{1}{\sqrt{2}}(1,-1)$ (stretch $2$, red). \textbf{Right:} For the $3\times3$ worked example, the plane eigenspace $E(1,A)=\spn(e_2,e_3)$ ($x+y+z=0$, blue) and the line eigenspace $E(4,A)=\spn(e_1)$ (red) are orthogonal complements in $\RR^3$.}
\end{fig}

\noindent\textbf{The question answered.} Yes: $A=\begin{psmallmatrix}3&1\\1&3\end{psmallmatrix}$ is symmetric, with orthonormal eigenvectors $u_1=\frac{1}{\sqrt{2}}(1,1)$ ($\lambda=4$) and $u_2=\frac{1}{\sqrt{2}}(1,-1)$ ($\lambda=2$). Writing $v=au_1+bu_2$ on the unit circle ($a^2+b^2=1$) gives $Av=4au_1+2bu_2$, which traces the ellipse with perpendicular semi-axes $4$ and $2$.

\noindent\textbf{Where this is used (Notes Ch.~20).} Transition operators of reversible Markov chains are self-adjoint under a stationary-distribution weighted inner product, yielding real eigenvalues whose second-largest modulus gives the spectral gap governing mixing speed. Adjacency matrices of undirected graphs and Hessian matrices of smooth functions are symmetric, explaining why their principal directions are mutually orthogonal.

\begin{shortcuts}
  \ii \textbf{$2\times2$ symmetric:} Once one eigenvector $(a,b)$ is found for distinct eigenvalues, the second is automatically $(-b,a)$ before normalizing (LADR~7.22). (If eigenvalues coincide, $A=\lambda I$.)
  \ii \textbf{Distinct eigenvalues:} Normalize eigenvectors directly; Gram--Schmidt is only needed inside eigenspaces of dimension $\ge2$.
  \ii \textbf{Remaining $2$-plane in $\RR^3$:} If eigenvalues are $\lambda$ (multiplicity $2$) and $\mu$ with unit eigenvector $e_1$ for $\mu$, then $E(\lambda,A)=(\spn e_1)^\perp$.
  \ii \textbf{Matrix powers:} $A^m = QD^m Q^{\mathsf T}$ and $p(A)=Qp(D)Q^{\mathsf T}$. For the question:
  \[
    A^m = \frac{4^m}{2}\begin{bmatrix}1&1\\1&1\end{bmatrix}+\frac{2^m}{2}\begin{bmatrix}1&-1\\-1&1\end{bmatrix}.
  \]
  \ii \textbf{Quick check:} Verify $Aq_j=\lambda_j q_j$ and $q_i\cdot q_j=\delta_{ij}$.
\end{shortcuts}

\begin{trapbox}
  \ii \textbf{Same eigenspace $\not\Rightarrow$ orthogonal:} $(1,-1,0)$ and $(1,0,-1)$ in $E(1,A)$ above have dot product $1$; orthogonalize within each eigenspace.
  \ii \textbf{Prop.~102(b) fails over $\CC$:} $T=iI$ on $\CC^n$ is orthogonally diagonalizable, yet $T^*=-iI\neq T$.
  \ii \textbf{Complex symmetric $\neq$ Hermitian:} Remember Ex.~74 ($M_2^{\mathsf T}=M_2$, eigenvalues $1\pm i\notin\RR$).
\end{trapbox}

\subsection{The complex spectral theorem (Ch.~20)}

\begin{qbox}
  The matrix $A=\begin{psmallmatrix}1&-1\\1&1\end{psmallmatrix}$ is normal ($AA^*=A^*A$) but not self-adjoint, so Prop.~102 does not apply. Does $\CC^2$ still have an \emph{orthonormal} basis of eigenvectors of $A$?
\end{qbox}

Throughout this subsection $V\neq\{0\}$ is a \textbf{complex} inner product space.

\begin{theorem*}[Prop.~103: Normal operators are diagonalizable]\target{prop-103}
  If $T\in\LL(V)$ ($V$ complex) is normal, then $T$ is diagonalizable.
\end{theorem*}

The notes give five independent proofs, each connecting normality to a different earlier week's perspective.

\begin{ideabox}
  \textbf{Proof~1 (minimal polynomial):} If $T$ failed to be diagonalizable, $m_T$ would have a repeated root $(x-\lambda)^2$ \ilink{T8}{(T8)}, yielding a nonzero $S=q(T)$ with $S^2=0$. Normality passes to polynomials in $T$, so Prop.~98 gives $\|S^*Sv\|=\|S^2v\|=0$, forcing $\|Sv\|^2=\langle S^*Sv,v\rangle=0$, so $q(T)=0$, contradicting minimality of $m_T$.
\end{ideabox}
\begin{proof}[Proof 1 (Minimal polynomial)]
  \begin{steps}
    \ii Suppose $m_T$ has a repeated root: $m_T=(x-\lambda)^2 r$. Put $q=(x-\lambda)r$, so $\deg q < \deg m_T$ forces $S=q(T)\neq0$, while $m_T\mid q^2$ gives $S^2=0$ \ilink{T8}{(T8)}.
    \ii Since $T$ commutes with $T^*$, every polynomial $q(T)$ commutes with its adjoint $q(T)^*=\overline{q}(T^*)$ (LADR~7.5), so $S$ is normal.
    \ii By Prop.~98, $\|S^*Sv\|=\|S(Sv)\|=0$ for every $v\in V$, so $S^*S=0$.
    \ii Therefore $\|Sv\|^2=\langle S^*Sv,v\rangle=0$ for all $v\in V$, forcing $q(T)=S=0$, contradicting Step~1.
    \ii Hence $m_T$ has only distinct roots over $\CC$ (T9/FTA), so $T$ is diagonalizable \ilink{T8}{(T8)}.
  \end{steps}
\end{proof}

\begin{ideabox}
  \textbf{Proof~2 (real and imaginary parts):} Write $T=R+iS$ as in LADR~7.23, where $R,S$ are commuting self-adjoint operators. Each is diagonalizable by Prop.~102; because $RS=SR$, each eigenspace of $S$ is $R$-invariant, so we can diagonalize $R$ separately inside each eigenspace of $S$ to obtain common eigenvectors.
\end{ideabox}
\begin{proof}[Proof 2 (Cartesian decomposition, LADR~7.23)]
  \begin{steps}
    \ii Write $T=R+iS$ with $R,S$ self-adjoint and $RS=SR$ (LADR~7.23). By Prop.~102, both $R$ and $S$ are diagonalizable.
    \ii For any eigenvalue $\mu$ of $S$, if $v\in E(\mu,S)$ then $S(Rv)=RSv=R(\mu v)=\mu(Rv)$, so $E(\mu,S)$ is $R$-invariant.
    \ii Thus $R|_{E(\mu,S)}$ is diagonalizable \ilink{T8}{(T8)}, giving a basis of $E(\mu,S)$ consisting of eigenvectors of both $R$ and $S$.
    \ii Since $V=\bigoplus_\mu E(\mu,S)$ \ilink{T8}{(T8)}, combining these bases produces a basis of $V$ \ilink{T15}{(T15)} of vectors $v$ with $Rv=\rho v$, $Sv=\mu v$, hence $Tv=(\rho+i\mu)v$.
  \end{steps}
\end{proof}

\begin{ideabox}
  \textbf{Proof~3 (induction on $\dim V$):} Split $V$ into the eigenspace $\nul(T-\lambda I)$ and its orthogonal complement $\range(T-\lambda I)$ (Prop.~99). Normality keeps $\range(T-\lambda I)$ invariant under both $T$ and $T^*$, so induction applies to that smaller subspace.
\end{ideabox}
\begin{proof}[Proof 3 (Induction on $n=\dim V$)]
  \begin{steps}
    \ii Base $n=1$ is clear. Let $n>1$ and pick an eigenvalue $\lambda$ of $T$ \ilink{T7}{(T7)}. Put $U=\range(T-\lambda I)$. If $U=\{0\}$, $T=\lambda I$ and every basis is an eigenbasis; else $1\le \dim U < n$ \ilink{T13}{(T13)}.
    \ii $U$ is invariant under $T$ and $T^*$: for $u=(T-\lambda I)v\in U$, $Tu=(T-\lambda I)Tv\in U$ and $T^*u=(T-\lambda I)T^*v\in U$ (since $T^*$ commutes with $T$).
    \ii Thus $(T|_U)^*=T^*|_U$, so $T|_U$ is normal; by induction $U$ has an eigenbasis for $T$.
    \ii By Prop.~94(i) and Prop.~99, $U^\perp = \nul\bigl((T-\lambda I)^*\bigr) = \nul(T-\lambda I)$, every nonzero vector of which is a $\lambda$-eigenvector.
    \ii Combining bases across $V=U\oplus U^\perp$ (T5, T15) yields an eigenbasis of $V$.
  \end{steps}
\end{proof}

\begin{ideabox}
  \textbf{Proof~4 (Jordan chains):} Over $\CC$, $V$ has a basis of Jordan chains \ilink{T11}{(T11)}. If a chain $v_1,v_2,\dots,v_k$ had length $k\ge2$, then $v_1=(T-\lambda I)v_2$ would sit simultaneously in $\range(T-\lambda I)$ and $\nul(T-\lambda I)$, which are orthogonal by Prop.~99. A vector orthogonal to itself is $0$, impossible in a basis.
\end{ideabox}
\begin{proof}[Proof 4 (Jordan chains)]
  \begin{steps}
    \ii Suppose a Jordan chain $v_1,\dots,v_k$ for $\lambda$ had $k\ge2$ \ilink{T11}{(T11)}. Then $v_1=(T-\lambda I)v_2\in\range(T-\lambda I)$ and $v_1\in\nul(T-\lambda I)$.
    \ii Since $T-\lambda I$ is normal (Prop.~99(b)), Prop.~99(a) gives $\range(T-\lambda I)=\nul(T-\lambda I)^\perp$, so $v_1\perp v_1$, forcing $v_1=0$, contradiction.
    \ii Every Jordan chain has length $1$, so $T$ is diagonalizable \ilink{T11}{(T11)}.
  \end{steps}
\end{proof}

To prepare Proof~5---which automatically yields an \emph{orthonormal} eigenbasis---we use Gram--Schmidt to upgrade upper-triangular bases \ilink{T10}{(T10)} to orthonormal ones.

\begin{proposition*}[LADR 6.38: Schur's theorem]\target{name-schur}
  Every operator $T\in\LL(V)$ on a finite-dimensional complex inner product space has an upper-triangular matrix in some orthonormal basis of $V$.
\end{proposition*}
\begin{proof}
  By \ilink{T10}{(T10)}, choose a basis $v_1,\dots,v_n$ in which each $\spn(v_1,\dots,v_k)$ is $T$-invariant. Applying Gram--Schmidt \ilink{T3}{(T3)} yields an orthonormal basis $e_1,\dots,e_n$ with $\spn(e_1,\dots,e_k)=\spn(v_1,\dots,v_k)$ for all $k$, so the same spans are $T$-invariant, and $\MM(T)$ is upper-triangular \ilink{T10}{(T10)}.
\end{proof}

\begin{ideabox}
  \textbf{Proof~5 (Schur):} Put $\MM(T)=(a_{ij})$ in upper-triangular Schur form in an orthonormal basis $e_1,\dots,e_n$. Then $\|Te_1\|^2=|a_{11}|^2$ (first column) while $\|T^*e_1\|^2=|a_{11}|^2+|a_{12}|^2+\dots+|a_{1n}|^2$ (conjugate-transpose gives the first row). By Prop.~98, $\|Te_1\|=\|T^*e_1\|$, which forces $a_{12}=\dots=a_{1n}=0$. Proceeding row by row eliminates every off-diagonal entry.
\end{ideabox}
\begin{proof}[Proof 5 (Schur's theorem)]
  \begin{steps}
    \ii Choose an orthonormal basis $e_1,\dots,e_n$ where $\MM(T)=(a_{ij})$ is upper-triangular (Schur's theorem). By \ilink{T3}{(T3)} and Prop.~93, $\|Te_k\|^2=\sum_{i=1}^n |a_{ik}|^2$ (column $k$) and $\|T^*e_k\|^2=\sum_{j=1}^n |a_{kj}|^2$ (row $k$).
    \ii For $k=1$, $a_{i1}=0$ for $i>1$, so Prop.~98 ($\|Te_1\|^2=\|T^*e_1\|^2$) forces $|a_{11}|^2 = \sum_{j=1}^n |a_{1j}|^2$, hence $a_{12}=\dots=a_{1n}=0$.
    \ii Inductively, if rows $1,\dots,k-1$ vanish off the diagonal, then column $k$ also has $a_{ik}=0$ for $i<k$ (and $a_{ik}=0$ for $i>k$ by upper-triangularity), leaving only $a_{kk}$. Then Prop.~98 gives $|a_{kk}|^2 = \sum_{j=k}^n |a_{kj}|^2$, forcing $a_{k,k+1}=\dots=a_{kn}=0$.
    \ii Hence $\MM(T)$ is diagonal, and $e_1,\dots,e_n$ is an orthonormal eigenbasis.
  \end{steps}
\end{proof}

\begin{theorem*}[Prop.~104 / LADR 7.31: The complex spectral theorem]\target{prop-104}\target{name-spectral}
  Let $V$ be a complex inner product space and $T\in\LL(V)$. Then $T$ is normal if and only if $T$ is orthogonally diagonalizable.
\end{theorem*}
\begin{proof}
  \begin{steps}
    \ii ($\Rightarrow$) Already proved by Proof~5 above. Alternatively, using Prop.~103: $T$ is diagonalizable so $V=\bigoplus_i E(\lambda_i,T)$ \ilink{T8}{(T8)}; distinct eigenspaces are orthogonal by LADR~7.22; choosing an orthonormal basis in each $E(\lambda_i,T)$ \ilink{T3}{(T3)} yields an orthonormal eigenbasis of $V$ \ilink{T15}{(T15)}.
    \ii ($\Leftarrow$) Suppose $e_1,\dots,e_n$ is an orthonormal basis with $Te_i=\lambda_i e_i$. In this basis $\MM(T)=\operatorname{diag}(\lambda_1,\dots,\lambda_n)$, so by Prop.~93 $\MM(T^*)=\operatorname{diag}(\overline{\lambda_1},\dots,\overline{\lambda_n})$. Thus $T^*Te_i = |\lambda_i|^2 e_i = TT^*e_i$ for all $i$, forcing $T^*T=TT^*$.
  \end{steps}
\end{proof}

\noindent\textbf{Matrix form ($A=UDU^*$).} A complex square matrix $U$ with orthonormal columns satisfies $U^*U=I$ and $U^{-1}=U^*$; such a matrix is called \textbf{unitary}. By Prop.~104, a complex matrix $A$ is normal $\iff A=UDU^*$ with $U$ unitary and $D$ diagonal.

\begin{example*}[Worked example: the matrix of the question]
  We saw $A=\begin{psmallmatrix}1&-1\\1&1\end{psmallmatrix}$ satisfies $AA^*=2I=A^*A$. By \ilink{T12}{(T12)}, $\lambda^2-2\lambda+2=0$ gives eigenvalues $1\pm i$.
  For $\lambda=1+i$, $A-(1+i)I=\begin{psmallmatrix}-i&-1\\1&-i\end{psmallmatrix}$ yields eigenvector $(1,-i)$.
  For $\lambda=1-i$, $A-(1-i)I=\begin{psmallmatrix}i&-1\\1&i\end{psmallmatrix}$ yields eigenvector $(1,i)$.
  Notice $\langle(1,-i),(1,i)\rangle = 1\cdot 1 + (-i)\cdot\overline{i} = 1+(-i)(-i) = 0$ (orthogonal by LADR~7.22). Normalizing by $\sqrt{2}$ gives $A=UDU^*$:
  \[
    U = \frac{1}{\sqrt{2}}\begin{bmatrix}1&1\\-i&i\end{bmatrix}, \qquad
    D = \begin{bmatrix}1+i&0\\0&1-i\end{bmatrix}.
  \]
\end{example*}

\noindent\textbf{The question answered.} Yes: $A$ is normal, so Prop.~104 guarantees the orthonormal eigenbasis $\frac{1}{\sqrt{2}}(1,-i)$, $\frac{1}{\sqrt{2}}(1,i)$ of $\CC^2$. (Over $\RR$, $A$ has no real eigenvalue and hence no real eigenbasis.)

\begin{shortcuts}
  \ii \textbf{Unitary diagonalizability test over $\CC$:} Check $AA^*=A^*A$; no characteristic polynomial is needed (Prop.~104).
  \ii \textbf{Second eigenvector of a $2\times2$ normal matrix:} If eigenvalues are distinct, take any vector orthogonal to $v_1$ (remembering to conjugate in $\CC^2$).
  \ii \textbf{Polynomials in normal operators:} If $T$ is normal, $p(T)$ is diagonal in the same orthonormal eigenbasis with eigenvalues $p(\lambda_j)$, hence $p(T)$ is normal; and $p(T)$ is self-adjoint $\iff$ all $p(\lambda_j)\in\RR$.
  \ii \textbf{Unitary inverses:} $U^{-1}=U^*$ without row reduction.
\end{shortcuts}

\begin{trapbox}
  \ii \textbf{Diagonalizable $\not\Rightarrow$ orthogonally diagonalizable:} $\begin{psmallmatrix}1&1\\0&2\end{psmallmatrix}$ has distinct eigenvalues $1,2$, yet its eigenvectors $(1,0)$ and $(1,1)$ are not orthogonal (it is not normal).
  \ii \textbf{Over $\RR$, normal $\not\Rightarrow$ diagonalizable:} A $90^\circ$ rotation on $\RR^2$ is normal ($RR^{\mathsf T}=I$), yet has no real eigenvalues.
  \ii \textbf{Complex dot product needs conjugation:} $\langle(1,-i),(1,i)\rangle = 1\cdot\overline{1}+(-i)\overline{i}=0$; without conjugating the second vector, $1+(-i)i=2\neq0$.
\end{trapbox}

\section{Problem set}

\begin{prob}\target{prob-HW11-P1}
  \textbf{HW11 P1.}\ Let $T\colon\FF^n\to\FF^n$ be $T(z_1,\dots,z_n)=(0,z_1,\dots,z_{n-1})$, where $\FF$ is $\RR$ or $\CC$. Find $T^*$.
  \par\smallskip\noindent{\color{OliveGreen!75!black}\textbf{Hint:}} Use \ilink{prop-93}{Prop.~93} in the standard (orthonormal) basis.
\end{prob}

\begin{prob}\target{prob-HW11-P2}
  \textbf{HW11 P2.}\ Let $T\colon V\to W$ be linear between finite-dimensional inner product spaces. Prove that $T$ is injective if and only if $T^*$ is surjective.
  \par\smallskip\noindent{\color{OliveGreen!75!black}\textbf{Hint:}} \ilink{prop-94}{Prop.~94(ii)} and $(\nul T)^\perp=V \iff \nul T=\{0\}$.
\end{prob}

\begin{prob}\target{prob-HW11-P3}
  \textbf{HW11 P3.}\ Let $T\colon\FF^2\to\FF^2$ be the quarter-turn $T(x,y)=(-y,x)$.
  \begin{enumerate}[(a)]
    \ii Find $T^*$.
    \ii Prove that $T$ is normal but not self-adjoint.
    \ii Find all eigenvalues of $T$ when $\FF=\RR$ and when $\FF=\CC$.
  \end{enumerate}
  \par\smallskip\noindent{\color{OliveGreen!75!black}\textbf{Hint:}} \ilink{prop-93}{Prop.~93}; compute $TT^*$ and $T^*T$; characteristic polynomial $\lambda^2+1$.
\end{prob}

\begin{prob}\target{prob-HW11-P4}
  \textbf{HW11 P4.}\ Let $P$ be a projector ($P^2=P$) on a finite-dimensional inner product space. Prove that the following are equivalent:
  \begin{enumerate}[(a)]
    \ii $P$ is self-adjoint;
    \ii $P$ is normal;
    \ii $P$ is the orthogonal projection onto $\range P$.
  \end{enumerate}
  \par\smallskip\noindent{\color{OliveGreen!75!black}\textbf{Hint:}} (b)$\Rightarrow$(c) by \ilink{prop-99}{Prop.~99}: $\nul P=(\range P)^\perp$; for (c)$\Rightarrow$(a) decompose $v=u+u'$ and $w=x+x'$ along $\range P\oplus(\range P)^\perp$.
\end{prob}

\begin{prob}\target{prob-HW11-P5}
  \textbf{HW11 P5.}\ Let $T$ be self-adjoint and $U$ a $T$-invariant subspace. Prove that $U^\perp$ is $T$-invariant.
  \par\smallskip\noindent{\color{OliveGreen!75!black}\textbf{Hint:}} \ilink{prop-101}{Prop.~101}, Step~1.
\end{prob}

\begin{prob}\target{prob-HW11-P6}
  \textbf{HW11 P6.}\ For the matrix $A=\begin{bmatrix}2&1\\1&2\end{bmatrix}$:
  \begin{enumerate}[(a)]
    \ii Find an orthonormal basis of $\RR^2$ of eigenvectors of $A$.
    \ii Find an orthogonal $Q$ and a diagonal $D$ with $A=QDQ^{\mathsf T}$.
    \ii Find a formula for $A^m$, $m\ge1$.
  \end{enumerate}
  \par\smallskip\noindent{\color{OliveGreen!75!black}\textbf{Hint:}} \ilink{prop-102}{Real spectral theorem} and the $2\times2$ symmetric shortcut; $A^m=QD^mQ^{\mathsf T}$.
\end{prob}

\begin{prob}\target{prob-HW11-P7}
  \textbf{HW11 P7 (optional).}\ Let $T$ be normal on a finite-dimensional complex inner product space and $p$ a polynomial. Use the complex spectral theorem to prove that $p(T)$ is normal.
  \par\smallskip\noindent{\color{OliveGreen!75!black}\textbf{Hint:}} \ilink{prop-104}{Prop.~104} both ways: in an orthonormal eigenbasis of $T$, $p(T)$ is diagonal.
\end{prob}

\begin{prob}\target{prob-O1}
  \textbf{O1 (adjoint for a weighted inner product).}\ On $\RR^2$ with $\langle x,y\rangle=2x_1y_1+x_2y_2$, let $T(x_1,x_2)=(x_2,0)$. Find $T^*$, and compare its matrix with the transpose of $\MM(T)$ in the standard basis.
  \par\smallskip\noindent{\color{OliveGreen!75!black}\textbf{Hint:}} Method of \ilink{ex-72}{Example~72}: write $T^*(e_1),T^*(e_2)$ with unknowns and enforce $\langle Tv,w\rangle=\langle v,T^*w\rangle$; or rescale $e_1$ to get an orthonormal basis and use \ilink{prop-93}{Prop.~93}.
\end{prob}

\begin{prob}\target{prob-O2}
  \textbf{O2 (eigenspaces of $T^*$).}\ Let $T\in\LL(V)$ and $\lambda\in\FF$. Prove that $\dim E(\overline{\lambda},T^*)=\dim E(\lambda,T)$.
  \par\smallskip\noindent{\color{OliveGreen!75!black}\textbf{Hint:}} \ilink{name-ladr-7-7}{LADR~7.7} for $T-\lambda I$, whose adjoint is $T^*-\overline{\lambda}I$.
\end{prob}

\begin{prob}\target{prob-O3}
  \textbf{O3 (a normal complex symmetric matrix).}\ For $M=\begin{bmatrix}1&i\\i&1\end{bmatrix}$ on $\CC^2$: is $M$ self-adjoint? normal? Find an orthonormal basis of eigenvectors, if there is one.
  \par\smallskip\noindent{\color{OliveGreen!75!black}\textbf{Hint:}} Compare $M$ with $M^*$, then compute $MM^*$ and $M^*M$; find eigenvalues with \ilink{T12}{(T12)}, and check orthogonality via \ilink{name-ladr-7-22}{LADR~7.22}.
\end{prob}

\begin{prob}\target{prob-O4}
  \textbf{O4 (real eigenvalues).}\ Let $T$ be normal on a complex $V$. Prove that $T$ is self-adjoint if and only if every eigenvalue of $T$ is real.
  \par\smallskip\noindent{\color{OliveGreen!75!black}\textbf{Hint:}} $(\Rightarrow)$ is \ilink{prop-95}{Prop.~95}; $(\Leftarrow)$ uses an orthonormal eigenbasis (\ilink{prop-104}{Prop.~104}) and \ilink{prop-93}{Prop.~93}.
\end{prob}

\begin{prob}\target{prob-O5}
  \textbf{O5 (LADR 7B).}\ Let $T$ be normal on a complex $V$ with $T^9=T^8$. Prove that $T$ is self-adjoint and $T^2=T$.
  \par\smallskip\noindent{\color{OliveGreen!75!black}\textbf{Hint:}} In an orthonormal eigenbasis (\ilink{prop-104}{Prop.~104}), $\lambda^9=\lambda^8$ for every eigenvalue, so $\lambda\in\{0,1\}$; then use \ilink{prob-O4}{O4}.
\end{prob}

\begin{prob}\target{prob-O6}
  \textbf{O6 (nilpotent and self-adjoint).}\ Let $T$ be self-adjoint with $T^2=0$. Prove that $T=0$.
  \par\smallskip\noindent{\color{OliveGreen!75!black}\textbf{Hint:}} $\|Tv\|^2=\langle T^2v,v\rangle$ when $T$ is self-adjoint.
\end{prob}

\end{document}
